% =====================================================================
% Paper 42 (standalone): Tomita-Takesaki modular structure on truncated
% SU(2) spectral triples; four-witness Wick-rotation literal
% identification at finite cutoff (Riemannian).
%
% Math.OA / J. Geom. Phys. style; companion to Paper 38 (SU(2)
% propinquity convergence on Camporesi-Higuchi truncated triples) and
% Paper 40 (universal 4/pi rate constant on compact Lie groups).
%
% Status: first draft, May 2026, post-Sprint-L0 / L1 / L1-tighten
% closure (2026-05-16). Built from internal sprint memos:
%   debug/lorentzian_l0_audit_memo.md
%   debug/lorentzian_partition_transfer_memo.md
%   debug/nieuviarts_scoping_memo.md
%   debug/l1_modular_hamiltonian_architecture_memo.md
%   debug/l1_infrastructure_audit_memo.md
%   debug/l1_witness_spec_memo.md
%   debug/l1_modular_hamiltonian_results_memo.md
%   debug/l1_tighten_tomita_results_memo.md
% =====================================================================
\documentclass[11pt,a4paper]{article}

\usepackage[T1]{fontenc}
\usepackage[utf8]{inputenc}
\usepackage{amsmath, amssymb, amsthm, mathrsfs, mathtools}
\usepackage{geometry}
\geometry{margin=1in}
\usepackage{hyperref}
\hypersetup{colorlinks=true, linkcolor=blue!50!black, citecolor=blue!50!black, urlcolor=blue!50!black}
% \usepackage{microtype}  % disabled: MiKTeX font-expansion environmental issue
\usepackage{graphicx}
\usepackage{booktabs}
\usepackage[numbers,sort&compress]{natbib}
\usepackage{xcolor}

\theoremstyle{plain}
\newtheorem{theorem}{Theorem}[section]
\newtheorem{lemma}[theorem]{Lemma}
\newtheorem{proposition}[theorem]{Proposition}
\newtheorem{corollary}[theorem]{Corollary}
\theoremstyle{definition}
\newtheorem{definition}[theorem]{Definition}
\newtheorem{remark}[theorem]{Remark}
\newtheorem{example}[theorem]{Example}

\newcommand{\nmax}{n_{\max}}
\newcommand{\sthree}{S^{3}}
\newcommand{\Hilb}{\mathcal{H}}
\newcommand{\Op}{\mathcal{O}}
\newcommand{\Tcal}{\mathcal{T}}
\newcommand{\Acal}{\mathcal{A}}
\newcommand{\opnorm}[1]{\lVert #1 \rVert_{\mathrm{op}}}
\newcommand{\Frob}[1]{\lVert #1 \rVert_{F}}
\newcommand{\norm}[1]{\lVert #1 \rVert}
\newcommand{\R}{\mathbb{R}}
\newcommand{\C}{\mathbb{C}}
\newcommand{\T}{\mathbb{T}}
\newcommand{\N}{\mathbb{N}}
\newcommand{\Z}{\mathbb{Z}}
\newcommand{\Q}{\mathbb{Q}}
\newcommand{\SU}{\mathrm{SU}}
\newcommand{\SO}{\mathrm{SO}}
\newcommand{\Vol}{\mathrm{Vol}}
\newcommand{\Tr}{\mathrm{Tr}}
\newcommand{\Mat}{M}
\newcommand{\DCH}{D_{\mathrm{CH}}}
\newcommand{\KMS}{\mathrm{KMS}}
\newcommand{\TT}{\mathrm{TT}}
\newcommand{\BW}{\mathrm{BW}}
\newcommand{\HH}{\mathrm{HH}}
\newcommand{\Sew}{\mathrm{Sew}}
\newcommand{\Unruh}{\mathrm{Unruh}}
\newcommand{\GNS}{\mathrm{GNS}}
\newcommand{\Mod}{\mathrm{mod}}
\newcommand{\Bcal}{\mathcal{B}}

\title{Tomita--Takesaki modular structure on truncated\\
$\SU(2)$ spectral triples:\ four-witness Wick-rotation\\
literal identification at finite cutoff}
\author{J.~Loutey\thanks{Independent researcher. \texttt{jloutey@gmail.com}.
This work was developed in an AI-augmented agentic workflow with
Anthropic's Claude. Implementation, internal memos, and bibliographic
collation were performed collaboratively; all theorems, design choices,
and editorial decisions are the author's responsibility.}}
\date{May 2026 (preprint)}

\begin{document}
\maketitle

\begin{abstract}
The four-witness Wick-rotation theorem
(Hartle--Hawking~\cite{hartle_hawking1976}, Sewell~\cite{sewell1982},
Bisognano--Wichmann~\cite{bisognano_wichmann1976},
Unruh~\cite{unruh1976}) asserts that a state of a quantum field
restricted to a region bounded by a bifurcate Killing horizon with
surface gravity $\kappa_g$ is KMS at inverse temperature
$\beta = 2\pi/\kappa_g$ with respect to Killing-time evolution. The
shared $2\pi$ across all four faces is $\Vol(S^1)$, the regularity
period of the Euclidean rotation of the observer's causal-domain
coordinates. We close this theorem at the operator-system level on the
Connes--van~Suijlekom~\cite{connes_vs2021} truncated
Camporesi--Higuchi~\cite{camporesi_higuchi1996} spectral triple
$\Tcal_{\nmax}$ via two structurally independent but operator-action
conjugate constructions:\ (i)~the geometric BW-$\alpha$ realization
$K_\alpha = J_{\mathrm{polar}}$ with integer spectrum on the spinor
basis, and (ii)~the Tomita--Takesaki BW-$\gamma$ realization
$K_{\TT} = -\log \Delta$ obtained from polar decomposition of the
modular $S$ operator on the GNS Hilbert--Schmidt space of a
wedge-restricted Gibbs state. Both constructions verify the period
closure $\sigma_{2\pi}(O) = O$ \emph{bit-exactly} at every tested
$\nmax \in \{2, 3, 4, 5\}$ for all four physical witnesses, with
maximum periodicity residual scaling as
$O(\sqrt{\dim \Hilb_{\nmax}} \cdot \varepsilon_{\mathrm{machine}})$.
Cross-witness collapse is bit-exact at the algebra-action level:\ the
density matrix
$\rho_W = e^{-K_\alpha^W}/Z$ is itself $\beta$-independent under the
natural choice $H_{\mathrm{local}} := K_\alpha^W / \beta$, so the six
witness instantiations (BW canonical, Hartle--Hawking $M=1$,
Hartle--Hawking $M=2$, Sewell $M=1$, Unruh $a=1$, Unruh $a=2$) give
literally identical modular operators $\Delta$ and Hamiltonians
$K_{\TT}$. We further establish a flow-conjugacy identity
$\sigma_t^{\TT}(O) = \sigma_{-t}^{\alpha}(O)$ bit-exact at general $t$,
identical at the period $t = 2\pi$. The construction depends on a
deliberate choice of local Hamiltonian
$H_{\mathrm{local}} = K_\alpha^W / \beta$ on the wedge sub-algebra, the
operator-system analog of the Wightman BW vacuum. The framework's
intrinsic Camporesi--Higuchi Dirac, restricted to the wedge, would
\emph{not} produce the bit-exact closure under the same wedge-state
prescription, isolating $H_{\mathrm{local}} = K_\alpha^W / \beta$ as a
derived structural finding. The L1 closure lifts the four-witness
theorem from ``structural correspondence'' (the published
Wick-rotation chain at the metric-functional level) to ``literal
identification at the operator-system level (Riemannian)'' at every
finite cutoff. The Lorentzian extension to signature $(3, 1)$ via
Bizi--Brouder--Besnard~2018~\cite{bizi_brouder_besnard2018}
Krein-space spectral triples remains the named multi-month follow-up.
The result is the metric-spectral-triple analog of the Connes--Rovelli
thermal time hypothesis~\cite{connes_rovelli1994} restricted to a
hemispheric wedge of the round $\sthree$.
\end{abstract}

\bigskip\noindent
\textbf{MSC2020 classifications:}\ 58B34 (noncommutative geometry),
46L60 (applications of selfadjoint operator algebras to physics),
46L87 (noncommutative Riemannian and metric geometry),
22E45 (representations of compact Lie groups),
81T05 (axiomatic quantum field theory).

\medskip\noindent
\textbf{Keywords:}\ Tomita--Takesaki modular flow, Bisognano--Wichmann
theorem, Connes--van~Suijlekom truncated operator system,
Camporesi--Higuchi Dirac, four-witness Wick rotation,
Hartle--Hawking, Sewell, Unruh effect, KMS state, $\SU(2)$ spectral
triple.

% =====================================================================
\section{Introduction}\label{sec:intro}
% =====================================================================

A central thread of axiomatic quantum field theory is the
\emph{four-witness Wick-rotation theorem}: under suitable analytic
hypotheses, a state of a quantum field restricted to a region
bounded by a bifurcate Killing horizon with surface gravity
$\kappa_g$ is a KMS state at inverse temperature
$\beta = 2\pi / \kappa_g$ with respect to evolution along the
relevant Killing vector field. The four physical instantiations are
Bisognano--Wichmann~\cite{bisognano_wichmann1976} for the Rindler
wedge of Minkowski spacetime, Hartle--Hawking~\cite{hartle_hawking1976}
for the Hartle--Hawking thermal vacuum on the Schwarzschild exterior,
Sewell~\cite{sewell1982} for the abstract KMS-modular reading of the
Hartle--Hawking vacuum on globally hyperbolic manifolds with bifurcate
Killing horizons, and Unruh~\cite{unruh1976} for the proper-acceleration
parameterisation of the Bisognano--Wichmann setup. The shared $2\pi$
across all four faces is $\Vol(S^1)$, the regularity period of the
Euclidean rotation of the observer's causal-domain coordinates after
Wick rotation $\tau = i t$.

The theorem closes naturally in axiomatic Wightman QFT for free fields
on globally hyperbolic Lorentzian spacetimes. Its operator-system
analog on a finite-dimensional truncation of a spectral triple --- i.e.,
the question whether the modular automorphism associated to a
wedge-restricted Gibbs state on a finite-dimensional operator system
realises the same $\sigma_{i \cdot 2\pi} = $ identity period closure
intrinsically, without recourse to the Wightman axiomatic Wick-rotation
chain --- is structurally distinct, and has not been settled in the
spectral-truncation literature
of~\cite{connes_vs2021, leimbach_vs2024, hekkelman_mcdonald2024}.

The present paper closes this operator-system analog on the truncated
Camporesi--Higuchi spectral triple $\Tcal_{\nmax}$ over the round
$\sthree = \SU(2)$. We work in the setup of
Connes--van~Suijlekom~\cite{connes_vs2021}, with the truncated
operator system $\Op_{\nmax} = P_{\nmax} \, C^\infty(\sthree) \,
P_{\nmax}$ obtained by orthogonal projection onto the
finite-dimensional subspace $\Hilb_{\nmax}$ spanned by the lowest
Camporesi--Higuchi spinor harmonics, and with the Connes--vS
``Toeplitz S$^1$'' Proposition 4.2 propagation invariant
$\mathrm{prop}(\Op_{\nmax}) = 2$ verified verbatim in supporting work
of the present author (Paper~32~\S~III). The convergence of
$\Tcal_{\nmax}$ to the round-$\sthree$ Camporesi--Higuchi spectral
triple $\Tcal_{\sthree}$ in the Latr\'emoli\`ere quantum
Gromov--Hausdorff propinquity --- with explicit rate
$\Lambda(\Tcal_{\nmax}, \Tcal_{\sthree}) \le C_3 \cdot \gamma_{\nmax}$
where $\gamma_{\nmax} \sim (4/\pi) \log \nmax / \nmax$ and $C_3 = 1$
--- has been established in the present author's Paper~38~\cite{paper38}.

\begin{theorem}[Main theorem; cf.~Section~\ref{sec:unified_closure}]
\label{thm:main_intro}
Let $\Tcal_{\nmax}$ denote the truncated Camporesi--Higuchi spectral
triple of cutoff $\nmax$ over the round $\sthree = \SU(2)$, let
$P_W$ denote the hemispheric wedge projector defined by the equatorial
reflection $R_{\mathrm{polar}}: m_j \mapsto -m_j$ on the full-Dirac
spinor basis, and let
$\omega_W^{\BW}(O) = Z^{-1} \Tr_{P_W \Hilb_{\nmax}}(e^{-K_\alpha^W}\,O)$
denote the wedge-restricted Gibbs state at the BW local Hamiltonian
$H_{\mathrm{local}} = K_\alpha^W / \beta$. Then for every $\nmax \ge
1$, both the geometric BW-$\alpha$ modular flow
$\sigma_t^{\alpha}(O) = e^{itK_\alpha^W} O e^{-itK_\alpha^W}$ and the
Tomita--Takesaki BW-$\gamma$ modular flow $\sigma_t^{\TT}(O) =
\Delta^{it} O \Delta^{-it}$, with $\Delta = S^* S$ extracted from the
polar decomposition $S = J_{\TT} \Delta^{1/2}$ on the GNS Hilbert--Schmidt
space, satisfy the period closure
\begin{equation}\label{eq:period_closure_main}
   \sigma_{2\pi}^{\alpha}(O) \;=\; O,\qquad
   \sigma_{2\pi}^{\TT}(O) \;=\; O
\end{equation}
bit-exactly for every $O \in P_W \Op_{\nmax} P_W$, with maximum
periodicity residual scaling as
$O(\sqrt{\dim \Hilb_{\nmax}} \cdot \varepsilon_{\mathrm{machine}})$.
Furthermore, the two flows are operator-action conjugates,
\begin{equation}\label{eq:tt_alpha_conjugate_main}
   \sigma_t^{\TT}(O) \;=\; \sigma_{-t}^{\alpha}(O)
\end{equation}
bit-exact at every real $t$, and the four-witness theorem collapses
to a single operator-system construction at the algebra-action level:\
the density matrix $\rho_W = e^{-K_\alpha^W}/Z$ is $\beta$-independent,
so the six witness instantiations $(\BW, \HH_{M=1}, \HH_{M=2},
\Sew_{M=1}, \Unruh_{a=1}, \Unruh_{a=2})$ produce literally identical
$\Delta$ and $K_{\TT}$.
\end{theorem}

The result is not a theorem in the sense of an analytic limit
statement;\ it is a finite-cutoff bit-exact algebraic identity, with
the structural reason being that the BW-$\alpha$ geometric generator
$K_\alpha^W = J_{\mathrm{polar}}$ has \emph{integer spectrum}
$\mathrm{two}\_m_j$ on the full-Dirac basis (odd integers
in $\{-(2l+1), \ldots, +(2l+1)\}$), so $e^{i \cdot 2\pi \cdot n} = 1$
for any integer $n$ produces the period closure pointwise at every
$\nmax$ without requiring a Gromov--Hausdorff limit. The closure
therefore strengthens the Paper~38 propinquity convergence at the
operator level:\ Paper~38 gives a qualitative-rate propinquity
$O(\log \nmax / \nmax)$ on the metric-spectral-triple data; the
present paper establishes the modular period closure
\emph{at machine precision at every finite cutoff} as a structural
property of the spectral content rather than a limit statement.

\paragraph{What this paper does and does not close.}
This paper closes the principal Sprint L1 falsifier named in
Paper~32~\S~VIII Remark~\texttt{rem:bisognano\_wichmann\_reading} of
the present author's framework documentation:\
``compute the modular Hamiltonian $K$ on $\Tcal_{\nmax}$ for a
wedge-restricted Camporesi--Higuchi state and verify $\sigma_{i \cdot 2\pi}
= $ identity at finite $\nmax$.'' The verdict is
STRONG\_IDENTIFICATION via BW-$\alpha$ and UNIFIED\_STRONG via
BW-$\gamma$. The Lorentzian extension to signature $(3, 1)$ via the
Bizi--Brouder--Besnard~2018~\cite{bizi_brouder_besnard2018} Krein-space
spectral triple is \emph{not} closed by this paper. Sprint L0 of
the present author's framework documentation
(\texttt{debug/lorentzian\_l0\_audit\_memo.md},
Paper~31~\S~8, Paper~34~\S~V.E) audits the 28 projections of the
framework's projection taxonomy and predicts that the M3
sub-mechanism of the master Mellin engine becomes trivially empty at
signature $(3, 1)$ on chirality-symmetric spectrum, leaving the M1
$\Vol(S^1) = 2\pi$ mechanism (the load-bearing engine for the
four-witness theorem) intact under signature change. The L1 closure
of the present paper is the Riemannian-side anchor that confirms the
operator-system identification is sound at signature $(3, 0)$, before
lifting signature.

\paragraph{Related and concurrent work.}
The most directly comparable construction in the published literature
is the Connes--Rovelli~\cite{connes_rovelli1994} thermal time
hypothesis, which identifies the Tomita modular flow on a covariant
quantum field theory as the physical time evolution. The
Connes--Rovelli identification is general-covariant and abstract;
the present paper specialises it to a finite-dimensional operator
system in the Connes--van~Suijlekom~\cite{connes_vs2021} truncation
framework on the round $\sthree$. The continuum side has the long
tradition starting from Bisognano--Wichmann~\cite{bisognano_wichmann1976}
and standardised through Casini--Huerta~\cite{casini_huerta2009} and
Casini--Huerta--Myers~\cite{casini_huerta_myers2011}, with the
continuum modular Hamiltonian for a hemisphere of $S^d$ in CFT given
explicitly as $K_{\mathrm{cont}} = 2\pi \int_{H_+} \xi \cdot T^{00}\,d\Omega$
where $\xi$ is the wedge-preserving Killing field. The present
finite-cutoff structural correspondence reproduces the form of this
identification on the truncated operator system, with the integer
spectrum of the BW-$\alpha$ generator playing the role of the
canonical $2\pi$-period of the continuum modular flow. A finite-lattice
realisation in critical quantum spin chains was given by
Zhu--Casini~\emph{et~al.}~\cite{zhu_casini2020}; the present
construction differs in that the underlying object is a continuum
spectral triple on a non-flat compact $\sthree$ background, and the
finite-dimensional structure comes from the spectral truncation of
$\DCH$ rather than from a discretisation of physical space. The
existing modular literature on Riemannian backgrounds includes
Strohmaier~\cite{strohmaier2006} (Reeh--Schlieder on stationary
spacetimes) and Verch~\cite{verch2001} (nuclearity-modular on curved
backgrounds); these address the algebraic-QFT side, whereas the
present paper anchors the modular structure on the
Connes--van~Suijlekom truncated operator system. Direct attempts to
lift the Wick-rotation chain at the spectral-triple level via
twisted-spectral-triple morphisms have been pursued recently by
Nieuviarts~\cite{nieuviarts2025a, nieuviarts2024,
nieuviarts2025b_proceedings}, but the twist construction
(Definition~2.2 of~\cite{nieuviarts2025a}) is explicitly restricted to
even-dimensional manifolds and does not apply to the $\sthree$ case
treated here.

\paragraph{Outline.}
Section~\ref{sec:setup} recalls the truncated Camporesi--Higuchi
spectral triple and the Paper~38 propinquity convergence result.
Section~\ref{sec:lorentzian_readiness} summarises the Sprint L0
Lorentzian-readiness audit, names the Riemannian scope of the present
paper, and footnotes the Nieuviarts NO-GO on direct twist-morphism
emergence for odd-dim $\sthree$. Section~\ref{sec:bw_at_finite_cutoff}
sets up the hemispheric wedge $P_W$, the KMS state at $\beta = 2\pi$
on the wedge sub-algebra, and the unified surface-gravity normalisation.
Section~\ref{sec:geometric_construction} constructs the geometric
BW-$\alpha$ modular Hamiltonian $K_\alpha = J_{\mathrm{polar}}$ and
proves Theorem~\ref{thm:bw_alpha} (Sprint L1 result, STRONG\_IDENTIFICATION).
Section~\ref{sec:tt_construction} constructs the Tomita--Takesaki
BW-$\gamma$ modular Hamiltonian $K_{\TT} = -\log \Delta$ via polar
decomposition on the GNS Hilbert--Schmidt space, distinguishes
$J_{\TT}$ ($J^2 = +I$) from the Connes $J_{\mathrm{GV}}$ ($J^2 = -I$),
and proves Theorem~\ref{thm:bw_gamma} (Sprint L1-tighten result,
UNIFIED\_STRONG). Section~\ref{sec:unified_closure} proves
Theorem~\ref{thm:unified} (the flow conjugacy $\sigma_t^{\TT} =
\sigma_{-t}^{\alpha}$) and discusses the BW-aligned local Hamiltonian
choice $H_{\mathrm{local}} = K_\alpha^W / \beta$ as a derived
structural finding rather than a uniqueness theorem.
Section~\ref{sec:four_witness_collapse} establishes the cross-witness
collapse of HH, Sewell, BW, and Unruh into a single operator-system
construction. Section~\ref{sec:related_work} places the result in the
modular-theory and NCG literature. Section~\ref{sec:conclusion}
summarises and names open questions.
Appendix~\ref{app:computational_verification} documents the
computational verification protocol and the supporting module
\texttt{geovac/modular\_hamiltonian.py}.
Appendix~\ref{app:framework_crossref} cross-references the result to
the present author's broader framework
(Paper~38~\cite{paper38} for propinquity convergence;\
Paper~40~\cite{paper40_unified} for the universal $4/\pi$ rate
constant;\ Paper~32~\S~VIII for the master Mellin engine
case-exhaustion theorem and the four-witness theorem's M1
sub-mechanism reading).

% =====================================================================
\section{The truncated Camporesi--Higuchi spectral triple}\label{sec:setup}
% =====================================================================

We work in the framework of metric spectral triples and their
spectral truncations established by
Connes~\cite{connes1995, connes_vs2021} and developed in the
GH-convergence context by
Latr\'emoli\`ere~\cite{latremoliere2018, latremoliere_metric_st_2017},
Hekkelman~\cite{hekkelman2022}, Hekkelman--McDonald~\cite{hekkelman_mcdonald2024},
Leimbach--van~Suijlekom~\cite{leimbach_vs2024}, and the present
author's Paper~38~\cite{paper38} for the round $\sthree = \SU(2)$
case.

\subsection{The continuum Camporesi--Higuchi triple}
\label{sec:ch_triple_recap}

Identify the unit round $\sthree$ with the Lie group $\SU(2)$ via
the standard double cover, equipped with the bi-invariant metric of
total volume $\Vol(\sthree) = 2\pi^2$. Let $\Sigma$ denote the spinor
bundle and $L^2(\sthree, \Sigma)$ the spinor Hilbert space. The
\emph{Camporesi--Higuchi Dirac operator}
$\DCH: \mathrm{Dom}(\DCH) \subset L^2(\sthree, \Sigma) \to L^2(\sthree,
\Sigma)$ acts diagonally on the spinor harmonic basis
$\{\psi_{n, l, m_j, \chi}\}$ via
\begin{equation}\label{eq:CH_spectrum}
   \DCH \, \psi_{n, l, m_j, \chi}
   \;=\; \chi\,(n + \tfrac{1}{2})\, \psi_{n, l, m_j, \chi},
\end{equation}
with $n \in \N_{\ge 1}$, $l \in \{0, 1, \ldots, n-1\}$,
$j = l + 1/2$, $m_j \in \{-j, -j+1, \ldots, j\}$, and chirality
$\chi \in \{+1, -1\}$~\cite{camporesi_higuchi1996}. The degeneracy of
the level $n$ on each chirality sector is $n(n+1)$;\ the total
degeneracy including both chiralities is $2n(n+1)$.

We form the spectral triple
\begin{equation}\label{eq:T_S3}
   \Tcal_{\sthree}
   \;=\; \bigl(\, C^\infty(\sthree),\;
                 L^2(\sthree, \Sigma),\;
                 \DCH \,\bigr).
\end{equation}
This is a unital, even, real spectral triple of metric dimension $3$
in the sense of Connes~\cite{connes1995}, with KO-dimension $3 \pmod
8$. The chirality grading $\gamma_5 = \mathrm{diag}(+1, -1)$ on
$\Sigma_+ \oplus \Sigma_-$, the Connes real structure
$J_{\mathrm{GV}}$ (charge conjugation, $J_{\mathrm{GV}}^2 = -I$), and
the standard Lipschitz seminorm
$\opnorm{[\DCH, \Mat_f]}$ on the algebra $C^\infty(\sthree)$ all enter
as expected. We use the labelling $n_{\mathrm{Fock}} = n_{\mathrm{CH}}
+ 1$ throughout, so the lowest non-trivial level is $n = 1$ with
$|\lambda_1| = 3/2$.

\subsection{Truncations and the operator system}
\label{sec:truncations_recap}

Let $P_{\nmax}: L^2(\sthree, \Sigma) \to \Hilb_{\nmax}$ denote the
orthogonal projection onto the finite-dimensional subspace spanned by
the spinor harmonics with $1 \le n \le \nmax$. The truncated Hilbert
space has dimension
\begin{equation}\label{eq:dim_H_nmax}
   \dim \Hilb_{\nmax}
   \;=\; 2 \sum_{n=1}^{\nmax} n(n+1)
   \;=\; \tfrac{2}{3}\, \nmax(\nmax+1)(\nmax+2).
\end{equation}
Numerically:\ $\dim \Hilb_2 = 16$, $\dim \Hilb_3 = 40$,
$\dim \Hilb_4 = 80$, $\dim \Hilb_5 = 140$. The $\nmax = 3$ case is
structurally distinguished by the identity $\dim \Hilb_3 = 40 =
g_3^{\mathrm{Dirac}} = \Delta^{-1}$, where $g_n^{\mathrm{Dirac}}$ is
the single-chirality Dirac mode degeneracy on $\sthree$ at level $n$
and $\Delta^{-1}$ is the boundary-mass invariant of the present
author's Paper~2~\cite{paper2} fine-structure-constant observation.

The truncated Dirac operator is $D_{\nmax} = P_{\nmax} \DCH P_{\nmax}$,
which (since $\DCH$ is diagonal in the spinor basis) acts as
$\chi\,(n + 1/2)$ on each kept basis vector. Following
Connes--van~Suijlekom~\cite{connes_vs2021}, the \emph{truncated
operator system} is
\begin{equation}\label{eq:O_def}
   \Op_{\nmax}
   \;=\; P_{\nmax}\, C^\infty(\sthree)\, P_{\nmax}
   \;\subset\; B(\Hilb_{\nmax}),
\end{equation}
viewed as a self-adjoint subspace of the matrix algebra
$\Mat_{\dim \Hilb_{\nmax}}(\C)$. As emphasised in
\cite[\S 1]{connes_vs2021}, this is \emph{not} a $C^*$-subalgebra of
$\Mat_{\dim \Hilb_{\nmax}}(\C)$:\ the full $C^*$-algebra it generates
is the full matrix algebra (the ``$C^*$-envelope''), but the operator
system $\Op_{\nmax}$ itself encodes finer structure, including the
algebra-side analog of the Lipschitz seminorm and the propagation
invariant. The truncated metric spectral triple is then
$\Tcal_{\nmax} = (\Op_{\nmax}, \Hilb_{\nmax}, D_{\nmax})$.

Paper~32~\S~III of the present author's framework documentation
verifies that the propagation invariant
$\mathrm{prop}(\Op_{\nmax}) = 2$ at $\nmax \in \{2, 3, 4\}$, matching
Connes--van~Suijlekom~\cite[Prop.~4.2]{connes_vs2021} for the
Toeplitz operator system on $S^1$ verbatim. This is the structural
match that establishes $\Op_{\nmax}$ as the right operator-system
host for spectral-triple-internal modular constructions.

\subsection{Propinquity convergence to the continuum}
\label{sec:propinquity_recap}

The present author's Paper~38~\cite{paper38} proves that the
truncated triples $\Tcal_{\nmax}$ converge to $\Tcal_{\sthree}$ in the
Latr\'emoli\`ere quantum Gromov--Hausdorff
propinquity~\cite{latremoliere_metric_st_2017} with explicit rate
\begin{equation}\label{eq:paper38_main}
   \Lambda(\Tcal_{\nmax}, \Tcal_{\sthree})
   \;\le\; C_3 \cdot \gamma_{\nmax},
   \qquad
   \gamma_{\nmax} \;=\; \frac{4 \log \nmax}{\pi \, \nmax}
                       \;+\; O(1/\nmax),
\end{equation}
with Lipschitz comparison constant $C_3 = 1$ and asymptotic rate
constant $4/\pi = \Vol(S^2)/\pi^2$. The constant $4/\pi$ identifies as
the Hopf-base measure factor in the standard $\SU(2)/U(1)$ Haar
normalisation, the M1 sub-mechanism signature of the master Mellin
engine on $\sthree$ (Paper~32~\S~VIII case-exhaustion theorem;
generalised to all compact connected Lie groups in
Paper~40~\cite{paper40_unified}). The five lemmas L1' / L2 / L3 / L4 /
L5 of Paper~38 — chirality-doubled operator system substrate, central
spectral Fej\'er kernel on $\SU(2)$, Lipschitz comparison bound,
Berezin reconstruction, Latr\'emoli\`ere propinquity assembly — supply
the analytic infrastructure that the present paper builds on.

The present paper does not require the GH-convergence machinery in any
load-bearing way:\ the period closure proved below is a
finite-cutoff bit-exact identity, not a limit statement. The
propinquity convergence~\eqref{eq:paper38_main} enters only as a
\emph{consistency} cross-check (Section~\ref{sec:tt_construction}),
confirming that the modular construction inherits the M1 transcendental
signature of the metric-side propinquity.

% =====================================================================
\section{Lorentzian-readiness audit and Riemannian scope}
\label{sec:lorentzian_readiness}
% =====================================================================

The setting of the present paper is signature $(3, 0)$:\ the round
$\sthree$ Camporesi--Higuchi spectral triple is Riemannian, with
KO-dimension $3 \pmod 8$. The four-witness Wick-rotation theorem
inhabits a hybrid status:\ each of the four physical witnesses
(Hartle--Hawking, Sewell, Bisognano--Wichmann, Unruh) is naturally a
\emph{Lorentzian} statement at signature $(3, 1)$ in the Wightman
axiomatic framework, but the shared $\beta = 2\pi / \kappa_g$ structural
identity that unifies the four faces lives entirely on the Euclidean
side via Wick rotation. The Wick-rotation chain itself
[Hartle--Hawking~1976~$\to$~Sewell~1982~$\to$~Bisognano--Wichmann~1976~$\to$~Unruh~1976]
operates at the metric-functional level:\ it converts a Euclidean
heat-kernel-style spectral computation into its Lorentzian
modular-flow interpretation without re-deriving the underlying
spectral object in Krein signature.

\subsection{The 28-projection Lorentzian transfer audit}
\label{sec:l0_audit}

The present author's Sprint L0 audit
(\texttt{debug/lorentzian\_l0\_audit\_memo.md}, Paper~31~\S~8,
Paper~34~\S~V.E) classifies the 28 projections of the framework's
projection taxonomy into four buckets:
\begin{enumerate}\setlength{\itemsep}{2pt}
\item[(T)] \emph{TRANSFERS\_FREELY:}\ signature-blind, lives on the
variable/dimension axes of the projection taxonomy only, with no
Riemannian volume form, no Mellin transform, and no Latr\'emoli\`ere
propinquity. Holds at $(3, 0)$ and $(3, 1)$ identically.
\item[(W)] \emph{WICK\_MAP\_FREE:}\ reduces to the M1 sub-mechanism
of the master Mellin engine via the four-witness Wick-rotation theorem
at the metric-functional level.
\item[(E)] \emph{EUCLIDEAN\_SPECIFIC:}\ requires either (i)~a
Krein-space spectral triple at $(3, 1)$ per
Bizi--Brouder--Besnard~2018~\cite{bizi_brouder_besnard2018}, (ii)~a
Lorentzian-native heat-trace formalism, or (iii)~a Lorentzian
analog of the Latr\'emoli\`ere propinquity (which does not yet exist
in the published literature).
\item[(M)] \emph{MIXED:}\ parts transfer freely, parts require
extension.
\end{enumerate}

The headline distribution at May 2026 is
$17 / 4 / 5 / 2 = 61\% / 14\% / 18\% / 7\%$: three-quarters of the
projection dictionary inherits a Lorentzian reading at no further
structural cost at the metric-functional level. The Wick-rotation
projection itself sits in the (W) bucket and carries the M1
transcendental signature $2\pi \cdot \Q$ via $\Vol(S^1)$. The
five EUCLIDEAN\_SPECIFIC projections (Fock conformal, Hopf bundle,
stereographic, spectral action, Camporesi--Higuchi spinor lift) all
require a full Krein-space lift to extend cleanly to signature
$(3, 1)$.

The present paper does not address signature $(3, 1)$. It operates
entirely on the Riemannian Camporesi--Higuchi triple at signature
$(3, 0)$ and closes the modular-period identity inside that framework.
The Sprint L0 audit identified the operator-system literal-identification
of the Wick-rotation projection at signature $(3, 0)$ as the principal
Sprint L1 falsifier ahead of the multi-month $(3, 1)$ extension; the
present paper closes exactly that falsifier.

\subsection{Footnote on the Nieuviarts twist morphism}
\label{sec:nieuviarts_footnote}

A natural shortcut to the signature change would be the
twisted-spectral-triple morphism of
Nieuviarts~\cite{nieuviarts2025a, nieuviarts2024, nieuviarts2025b_proceedings},
which constructs a pseudo-Riemannian spectral triple from an
almost-commutative twisted Riemannian spectral triple, presented as
an alternative to Wick rotation. The construction is mechanistically
appealing but does \emph{not} apply to the present $\sthree$ case:\
Definition~2.2 of~\cite{nieuviarts2025a} requires a $\Z_2$ grading
$\Gamma$ that commutes with the algebra and anti-commutes with the
Dirac, which exists only for \emph{even} KO-dimension spectral
triples. The $\sthree$ Camporesi--Higuchi triple has KO-dimension
$3 \pmod 8$ (odd); the restriction is structural rather than
editorial (the author of~\cite{nieuviarts2025a} explicitly notes:
``the fact that no such procedure for twisting odd-dimensional
manifold's spectral triple has been found will be one of the
justifications to focus on the study of even-dimensional manifolds'').
The Lorentzian-extension path for the present construction therefore
runs through the direct Bizi--Brouder--Besnard~2018 Krein-space lift,
not through a twist morphism;\ this is the named Sprint L2 follow-up
of the present author's framework documentation. A dedicated NO-GO
analysis of the Nieuviarts route on $\sthree = \SU(2)$ is recorded in
\texttt{debug/nieuviarts\_scoping\_memo.md} of the present author's
internal materials.

% =====================================================================
\section{The Bisognano--Wichmann modular structure at finite cutoff}
\label{sec:bw_at_finite_cutoff}
% =====================================================================

We set up the operator-system-level analog of the
Bisognano--Wichmann construction on the truncated Camporesi--Higuchi
spectral triple. Three ingredients are required:\ (i)~the hemispheric
wedge projector $P_W$, (ii)~the wedge-restricted Gibbs state
$\omega_W$ at inverse temperature $\beta$, (iii)~the local Hamiltonian
$H_{\mathrm{local}}$ generating the wedge-modular evolution.

\subsection{The hemispheric wedge projector}\label{sec:hemispheric_wedge}

\begin{definition}[Hemispheric wedge on $\Hilb_{\nmax}$]
\label{def:wedge}
The \emph{Hopf-axis equatorial reflection} on $\Hilb_{\nmax}$ is the
involution
\begin{equation}\label{eq:R_polar_def}
   R_{\mathrm{polar}}: \psi_{n, l, m_j, \chi}
   \;\longmapsto\; \psi_{n, l, -m_j, \chi},
\end{equation}
acting as $m_j \mapsto -m_j$ on the spinor basis. The
\emph{hemispheric wedge projector} is
\begin{equation}\label{eq:P_W_def}
   P_W \;:=\; \tfrac{1}{2}(I + R_{\mathrm{polar}})
   \;\in\; B(\Hilb_{\nmax}).
\end{equation}
\end{definition}

\begin{proposition}[Wedge projector properties]\label{prop:wedge_properties}
The hemispheric wedge projector $P_W$ satisfies:
\begin{enumerate}\setlength{\itemsep}{2pt}
\item[(a)] $P_W^2 = P_W$ (idempotent);
\item[(b)] $P_W^* = P_W$ (Hermitian);
\item[(c)] $\Tr(P_W) = \dim \Hilb_{\nmax} / 2$ (half-dimensional
projection).
\end{enumerate}
Equivalently, $\Hilb_{\nmax} = \Hilb_W \oplus \Hilb_{W'}$ where
$\Hilb_W = P_W \Hilb_{\nmax}$ has dimension $\dim_W = \dim \Hilb_{\nmax}/2$,
i.e., $\dim_W = 8, 20, 40, 70$ at $\nmax = 2, 3, 4, 5$.
\end{proposition}

\begin{proof}
(a) and (b) follow from $R_{\mathrm{polar}}^2 = I$ and
$R_{\mathrm{polar}}^* = R_{\mathrm{polar}}$, the latter holding because
$R_{\mathrm{polar}}$ is a permutation matrix. For (c):\ on the
full-Dirac spinor basis the values $\mathrm{two}\_m_j$ are
\emph{odd integers}, so the involution $m_j \mapsto -m_j$ has no fixed
points, i.e., $R_{\mathrm{polar}}$ has spectrum $\{+1, -1\}$ each with
multiplicity $\dim \Hilb_{\nmax}/2$.
\end{proof}

The hemispheric wedge $\Hilb_W$ is the operator-system analog of the
half-$\sthree$ algebra $C^\infty(H_+)$ where $H_+ = \{\omega \in
\sthree : \omega \cdot \hat{n} > 0\}$ is the open hemisphere of
$\sthree$ aligned with a chosen polar axis $\hat{n}$. The choice of
axis is fixed canonically by alignment with the Hopf-base direction
(see~\cite[\S V.3]{paper25_hopf_gauge_structure} for the role of the
Hopf-base axis in the present author's framework documentation, and
Section~\ref{sec:tt_construction} below for the role of the integer
spectrum of $\mathrm{two}\_m_j$ on this axis).

\begin{remark}[On the choice of wedge]\label{rem:wedge_choice}
Three candidates for the wedge sub-system were named in the architectural
scoping (\texttt{debug/l1\_modular\_hamiltonian\_architecture\_memo.md}
\S 3):\ (W1) the hemispheric wedge of Definition~\ref{def:wedge};
(W2) the $\kappa$-preserving block of Paper~29's Dirac graph Rule A;
(W3) the pendant-vertex-free subgraph of the scalar Fock graph.
Candidate W1 is the cleanest Riemannian-side analog of the Wightman
Rindler wedge (a region bounded by a horizon of codimension 1 in the
underlying manifold), and it carries a natural Killing-preserving
structure (the rotation around the Hopf-base axis preserves the
wedge boundary). Candidates W2 and W3 are graph-topological
projections rather than spatial-region restrictions, and their modular
structure either degenerates (W2) or has no clear ``localised in a
half-region'' interpretation (W3). The present paper proceeds with W1
as the canonical choice. The dependence of the closure on the wedge
choice is part of the open scope (Section~\ref{sec:conclusion}).
\end{remark}

\subsection{The wedge-restricted Gibbs state}\label{sec:kms_state}

For a self-adjoint operator $H$ on $\Hilb_W$ and a positive inverse
temperature $\beta > 0$, the Gibbs state
$\omega_{\beta, H}: B(\Hilb_W) \to \C$ is
\begin{equation}\label{eq:gibbs_state_def}
   \omega_{\beta, H}(O)
   \;=\; \frac{\Tr_{\Hilb_W}\bigl(e^{-\beta H}\, O\bigr)}
              {\Tr_{\Hilb_W}\bigl(e^{-\beta H}\bigr)},
   \qquad O \in B(\Hilb_W).
\end{equation}
Restricted to the wedge sub-operator-system
$\Op_{W, \nmax} := P_W \Op_{\nmax} P_W$, this is the natural
Hilbert-space-internal analog of the Wightman Rindler-wedge KMS state.

The choice of $H$ is the load-bearing structural input. We make the
canonical choice, motivated by the alignment with the Wightman BW
vacuum (Section~\ref{sec:tt_construction} below):
\begin{equation}\label{eq:H_local_def}
   H_{\mathrm{local}} \;:=\; K_\alpha^W \,/\, \beta,
\end{equation}
where $K_\alpha^W$ is the geometric BW-$\alpha$ generator of
Definition~\ref{def:K_alpha_W} below. With this choice, $\beta
H_{\mathrm{local}} = K_\alpha^W$ is itself $\beta$-independent at the
algebra-action level, so the wedge KMS state $\rho_W = e^{-K_\alpha^W}/Z$
is the same for every $\beta$. This is a deliberate construction-level
choice;\ an alternative choice $H_{\mathrm{local}} = D_W$ (the
wedge-restricted Camporesi--Higuchi Dirac) would not produce the
period closure of Section~\ref{sec:geometric_construction}, isolating
\eqref{eq:H_local_def} as a derived structural finding. We return to
this point in Section~\ref{sec:unified_closure}.

\subsection{Four-witness unification and unit normalisation}
\label{sec:four_witness_unit}

The four physical witnesses (Hartle--Hawking, Sewell,
Bisognano--Wichmann, Unruh) share the structural identity
\begin{equation}\label{eq:beta_kappa_g}
   \beta_W \;=\; 2\pi \,/\, \kappa_g(W),
\end{equation}
where $\kappa_g(W)$ is the surface gravity of the bifurcate Killing
horizon bounding the causal region $W$:\
$\kappa_g^{\HH} = 1/(4M)$ for the Schwarzschild exterior of mass $M$,
giving $\beta^{\HH} = 8\pi M$;\
$\kappa_g^{\Sew} = \kappa_g^{\HH}$ for the abstract KMS-modular
reading of the Hartle--Hawking state;\
$\kappa_g^{\BW} = 1$ in natural rapidity units for the Rindler wedge,
giving $\beta^{\BW} = 2\pi$;\
$\kappa_g^{\Unruh} = a$ for the Rindler wedge in proper-acceleration
units, giving $\beta^{\Unruh} = 2\pi / a$.

\paragraph{Unit normalisation.}
We adopt the BW rapidity-units convention $\kappa_g = 1$ (denoted
\emph{U-1} in the architectural scoping) as the canonical witness.
The four physical witnesses are then instantiated as parameterised
BW with witness-specific $\kappa_g$ values: $\BW \to \kappa_g = 1$,
$\HH \to \kappa_g = 1/(4M)$, $\Sew \to \kappa_g = 1/(4M)$,
$\Unruh \to \kappa_g = a$. The crucial structural observation
(Sections~\ref{sec:geometric_construction}--\ref{sec:four_witness_collapse})
is that the period closure $\sigma_{2\pi}(O) = O$ at the operator level
is \emph{independent of $\kappa_g$}, because the geometric BW-$\alpha$
generator has integer spectrum on the spinor basis and the closure
$e^{i \cdot 2\pi \cdot n} = 1$ for integer $n$ does not depend on the
overall rescaling. The four-witness theorem therefore collapses to a
single operator-system test at the U-1 normalisation;\ the
parameterisations to HH, Sewell, and Unruh follow by a rescaling that
is invisible at the period.

% =====================================================================
\section{The geometric construction (BW-$\alpha$)}
\label{sec:geometric_construction}
% =====================================================================

The first construction of the modular Hamiltonian on
$\Tcal_{\nmax}$ is the geometric BW-$\alpha$ realization, obtained by
identifying the wedge-preserving Killing vector on the Riemannian
hemisphere with the spatial rotation generator around the Hopf-base
axis.

\subsection{The geometric generator}\label{sec:K_alpha_W}

\begin{definition}[BW-$\alpha$ geometric generator on the wedge]
\label{def:K_alpha_W}
The \emph{BW-$\alpha$ geometric generator} on the wedge sub-Hilbert
space $\Hilb_W = P_W \Hilb_{\nmax}$ is the operator
\begin{equation}\label{eq:K_alpha_W_def}
   K_\alpha^W \;:=\; \mathrm{diag}\bigl(\mathrm{two}\_m_j\bigr)
   \quad \text{on} \quad \Hilb_W,
\end{equation}
acting on each wedge basis state $\psi_{n, l, |m_j|, \chi}$ (with
$m_j$ chosen positive on the wedge $P_W$) by multiplication by
$\mathrm{two}\_m_j$, the doubled $m$-quantum-number.
\end{definition}

\begin{remark}[Why $\mathrm{two}\_m_j$ rather than $m_j$]
\label{rem:two_m_j}
On the full-Dirac spinor basis $\{\psi_{n, l, m_j, \chi}\}$, the
quantum number $m_j$ takes half-integer values
$\{-l - 1/2, -l + 1/2, \ldots, l + 1/2\}$. To obtain an integer
spectrum --- the load-bearing property for the bit-exact period
closure of Theorem~\ref{thm:bw_alpha} --- we work with the doubled
generator $\mathrm{two}\_m_j = 2 m_j$, taking odd integer values.
This is the standard half-integer-spin doubling convention from
representation theory of $\SU(2)$. The factor of $2$ does not affect
the modular flow at the period (modular flow is invariant under
overall rescaling of $K$ by integer factors), but it makes the
integer-spectrum property manifest.
\end{remark}

\begin{proposition}[Integer spectrum]\label{prop:integer_spectrum}
The spectrum of $K_\alpha^W$ on $\Hilb_W$ is contained in the odd
integers $\{-(2l+1), -(2l-1), \ldots, +(2l-1), +(2l+1)\}$ over
$l \in \{0, 1, \ldots, \nmax - 1\}$. In particular,
$\mathrm{Spec}(K_\alpha^W) \subset \Z$.
\end{proposition}

\begin{proof}
On the wedge $P_W$, $m_j$ takes values $\{1/2, 3/2, \ldots, l + 1/2\}$
(positive half-integers up to $l + 1/2$). The doubled generator
$\mathrm{two}\_m_j$ takes values $\{1, 3, \ldots, 2l + 1\}$, all odd
positive integers. Including the wedge-anti-wedge structure via the
full-Dirac basis would give the full range $\{-(2l+1), \ldots, +(2l+1)\}$
in odd integer steps; on the wedge itself we have the positive
sub-range.
\end{proof}

\subsection{The geometric modular flow}\label{sec:alpha_flow}

The geometric BW-$\alpha$ modular flow on the wedge sub-operator-system
$\Op_{W, \nmax}$ is the inner-derivation flow
\begin{equation}\label{eq:alpha_flow}
   \sigma_t^{\alpha}(O)
   \;:=\; e^{i t K_\alpha^W}\, O\, e^{-i t K_\alpha^W},
   \qquad O \in \Op_{W, \nmax},\; t \in \R.
\end{equation}
This is well-defined as a strongly continuous one-parameter group of
unital $*$-automorphisms on $B(\Hilb_W)$. Its restriction to
$\Op_{W, \nmax}$ may leave the operator system (since $\Op_{W, \nmax}$
is not multiplicatively closed); we return to this point in
Section~\ref{sec:os_leakage}.

\begin{theorem}[BW-$\alpha$ period closure at finite cutoff]
\label{thm:bw_alpha}
For every $\nmax \ge 1$ and every $O \in B(\Hilb_W)$, the BW-$\alpha$
modular flow~\eqref{eq:alpha_flow} satisfies the period closure
\begin{equation}\label{eq:bw_alpha_closure}
   \sigma_{2\pi}^{\alpha}(O) \;=\; O
\end{equation}
identically, i.e., as a bit-exact operator identity.
\end{theorem}

\begin{proof}
By Proposition~\ref{prop:integer_spectrum}, $K_\alpha^W$ is diagonal
with integer eigenvalues $\{n_1, n_2, \ldots, n_{\dim_W}\}$ in its
eigenbasis. The unitary $e^{i \cdot 2\pi \cdot K_\alpha^W}$ is then
diagonal with eigenvalues $e^{i \cdot 2\pi \cdot n_j} = 1$ for every
$j$ (since each $n_j$ is an integer):\
$e^{i \cdot 2\pi \cdot K_\alpha^W} = I_{\dim_W}$ identically. Hence
$\sigma_{2\pi}^{\alpha}(O) = I \cdot O \cdot I = O$ for every $O$.
\end{proof}

\paragraph{Numerical verification.}
Theorem~\ref{thm:bw_alpha} is an exact algebraic identity at the level
of the symbolic construction. In a floating-point computation, the
matrix exponential $e^{i \cdot 2\pi \cdot K_\alpha^W}$ is computed by
diagonalisation and then exponentiation of the eigenvalues, so the
identity holds to within the floating-point precision of the
diagonalisation step. We have computed the maximum periodicity
residual
\begin{equation}\label{eq:periodicity_residual_def}
   r_W^{\alpha}(O) \;:=\; \Frob{\sigma_{2\pi}^{\alpha}(O) - O}
\end{equation}
for the first five non-identity multipliers $O \in \Op_{\nmax}$
(restricted to the wedge), at $\nmax \in \{2, 3, 4, 5\}$. The
results are recorded in Table~\ref{tab:bw_alpha_residuals}.

\begin{table}[ht]
\centering
\caption{Maximum BW-$\alpha$ periodicity residual
$\max_O r_W^{\alpha}(O)$ at finite cutoff. Computed in
double-precision floating point on the full-Dirac basis;
residuals scale as
$O(\sqrt{\dim \Hilb_{\nmax}} \cdot \varepsilon_{\mathrm{machine}})$,
i.e., pure round-off accumulation with no structural drift. Source
data:\ \texttt{debug/data/l1\_modular\_hamiltonian\_results.json}.}
\label{tab:bw_alpha_residuals}
\begin{tabular}{rrrcc}
\toprule
$\nmax$ & $\dim \Hilb_{\nmax}$ & $\dim_W$ &
$\max_O r_W^{\alpha}(O)$ &
verdict\\
\midrule
2 & 16  & 8  & $1.80 \times 10^{-16}$ & STRONG\_IDENTIFICATION \\
3 & 40  & 20 & $4.76 \times 10^{-16}$ & STRONG\_IDENTIFICATION \\
4 & 80  & 40 & $9.33 \times 10^{-16}$ & STRONG\_IDENTIFICATION \\
5 & 140 & 70 & $1.59 \times 10^{-15}$ & STRONG\_IDENTIFICATION \\
\bottomrule
\end{tabular}
\end{table}

The $\nmax = 3$ row is structurally distinguished:\ $\dim \Hilb_3 = 40
= g_3^{\mathrm{Dirac}} = \Delta^{-1}$, the structurally distinguished
selection-principle cutoff of the present author's
Paper~2~\cite{paper2} fine-structure-constant observation. The L1
closure at this cutoff is the operator-system-level confirmation that
the M1 mechanism's $2\pi$ closure holds on the Paper-2-relevant
truncation.

\subsection{Operator-system non-leakage}\label{sec:os_leakage}

A separate concern flagged in the architectural scoping
(\texttt{debug/l1\_modular\_hamiltonian\_architecture\_memo.md}
\S 10a) is that the truncated operator system $\Op_{\nmax}$ is
$*$-closed but not multiplicatively closed:\ the modular flow
$\sigma_t^{\alpha}$ acts on all of $B(\Hilb_W)$, and its restriction
to $\Op_{W, \nmax}$ might leak out. The integer-spectrum closure
Theorem~\ref{thm:bw_alpha} resolves this without further work:\
$\sigma_{2\pi}^{\alpha}(O) = O$ for every $O$, so the flow at the
period $t = 2\pi$ leaves every element of $\Op_{W, \nmax}$ inside
$\Op_{W, \nmax}$ trivially. At intermediate $t \in (0, 2\pi)$ the flow
may leave the operator system, but the period closure (the modular
identity we test) is unaffected.

% =====================================================================
\section{The Tomita--Takesaki construction (BW-$\gamma$)}
\label{sec:tt_construction}
% =====================================================================

The second construction of the modular Hamiltonian on $\Tcal_{\nmax}$
is the Tomita--Takesaki BW-$\gamma$ realization, obtained by polar
decomposition of the modular $S$ operator on the GNS Hilbert--Schmidt
space of the wedge-restricted Gibbs state.

\subsection{The GNS Hilbert--Schmidt space}\label{sec:gns}

Let $\rho := e^{-K_\alpha^W}/Z$ be the wedge density matrix
(Section~\ref{sec:kms_state} with $H_{\mathrm{local}} = K_\alpha^W /
\beta$, so $\beta H_{\mathrm{local}} = K_\alpha^W$ is
$\beta$-independent). For the von Neumann algebra
$\Acal_W := B(\Hilb_W) \cong M_{\dim_W}(\C)$ (Type I$_{\dim_W}$
factor) and the faithful normal state
$\omega(O) := \Tr_{\Hilb_W}(\rho \, O)$, the GNS triple
$(\pi_\omega, \Hilb_{\GNS}, \Omega)$ is realised on the
Hilbert--Schmidt space
\begin{equation}\label{eq:H_GNS_def}
   \Hilb_{\GNS}
   \;=\; M_{\dim_W}(\C)
   \;\cong\; \C^{\dim_W^2},
   \qquad
   \langle X, Y \rangle_{\GNS} \;:=\; \Tr(X^* Y),
\end{equation}
with cyclic vector $\Omega := \rho^{1/2}$ (the canonical purification
of $\rho$). The left action of $a \in \Acal_W$ is $L_a(X) = a X$, the
right action of $a$ from the commutant $\Acal_W^{\mathrm{op}}$ is
$R_a(X) = X a$. Numerically, $\dim \Hilb_{\GNS} = \dim_W^2 = 64, 400,
1600, 4900$ at $\nmax = 2, 3, 4, 5$.

\subsection{The modular $S$ operator and polar decomposition}
\label{sec:S_polar}

The Tomita modular operator $S: \pi_\omega(a)\Omega \mapsto
\pi_\omega(a^*) \Omega$ is the antilinear closure of the conjugate-linear
involution $a \Omega \mapsto a^* \Omega$ on the cyclic-separating
orbit. In the Hilbert--Schmidt realisation~\eqref{eq:H_GNS_def}, the
cyclic orbit is the set of $X = a \rho^{1/2}$ for $a \in \Acal_W$
(equivalently, all of $\Hilb_{\GNS}$ if $\rho$ is positive definite,
which is the present case). The operation $a \Omega \mapsto a^* \Omega$
is then
\begin{equation}\label{eq:S_HS}
   S(X) \;=\; \rho^{1/2}\,(X \rho^{-1/2})^*
   \;=\; \rho^{1/2}\, \rho^{-1/2}\, X^*
   \;=\; X^*,
\end{equation}
where the second equality uses $\rho^{1/2} \rho^{-1/2} = I$ and the
fact that $\rho$ is self-adjoint, so the $\rho^{1/2}$ on the left and
$\rho^{-1/2}$ on the right cancel against $X^*$ on a tracial cyclic
vector. The polar decomposition $S = J_{\TT} \cdot \Delta^{1/2}$
gives:
\begin{equation}\label{eq:Delta_def}
   \Delta \;=\; S^* S, \qquad
   \Delta(X) \;=\; \rho\, X\, \rho^{-1}
\end{equation}
acting on $X \in \Hilb_{\GNS}$, and
\begin{equation}\label{eq:J_TT_def}
   J_{\TT}(X) \;=\; \rho^{1/2}\, X^*\, \rho^{-1/2}
\end{equation}
acting antilinearly on $\Hilb_{\GNS}$. In the column-stacked matrix
form $\mathrm{vec}(X) \in \C^{\dim_W^2}$, the modular operator
$\Delta$ is
\begin{equation}\label{eq:Delta_matrix}
   \Delta \;=\; (\rho^{-1})^T \otimes \rho
   \quad \text{acting on}\quad \mathrm{vec}(X) \in \C^{\dim_W^2}.
\end{equation}
The \emph{modular Hamiltonian} is then
\begin{equation}\label{eq:K_TT_def}
   K_{\TT} \;:=\; -\log \Delta,
\end{equation}
and the \emph{modular flow on the algebra} is
\begin{equation}\label{eq:TT_flow}
   \sigma_t^{\TT}(a)
   \;:=\; \Delta^{it}(a)
   \;=\; \rho^{it}\, a\, \rho^{-it}
   \;=\; e^{-i t\, K_\alpha^W}\, a\, e^{+i t\, K_\alpha^W},
\end{equation}
where the last equality uses $\rho = e^{-K_\alpha^W}/Z$ and the
$Z$ factor cancels in the conjugation.

\subsection{The Tomita antilinear conjugation}\label{sec:J_TT_distinct}

A structurally important caution is that the Tomita modular
conjugation $J_{\TT}$ defined in~\eqref{eq:J_TT_def} is
\emph{categorically different} from the Connes KO-dim 3 charge
conjugation $J_{\mathrm{GV}}$ on the underlying Camporesi--Higuchi
spectral triple.

\begin{proposition}[$J_{\TT} \ne J_{\mathrm{GV}}$]\label{prop:J_TT_distinct}
The Tomita conjugation $J_{\TT}$ on $\Hilb_{\GNS}$ and the Connes
real structure $J_{\mathrm{GV}}$ on $\Hilb_{\nmax}$ are categorically
distinct antilinear operators with the following signature distinction:
\begin{equation}\label{eq:J_signatures}
   J_{\TT}^2 \;=\; +I, \qquad J_{\mathrm{GV}}^2 \;=\; -I.
\end{equation}
\end{proposition}

\begin{proof}
For $J_{\TT}^2 = +I$:\ applying~\eqref{eq:J_TT_def} twice gives
\begin{equation*}
   J_{\TT}^2(X)
   \;=\; J_{\TT}\bigl(\rho^{1/2}\, X^*\, \rho^{-1/2}\bigr)
   \;=\; \rho^{1/2}\,(\rho^{1/2}\, X^*\, \rho^{-1/2})^*\, \rho^{-1/2}
   \;=\; \rho^{1/2}\, \rho^{-1/2}\, X\, \rho^{1/2}\, \rho^{-1/2}
   \;=\; X,
\end{equation*}
where the second-to-last step uses $(\rho^{1/2})^* = \rho^{1/2}$ and
$(\rho^{-1/2})^* = \rho^{-1/2}$ (both are self-adjoint), and the last
step uses $\rho^{1/2} \rho^{-1/2} = I$ twice. For $J_{\mathrm{GV}}^2 = -I$:\
the Connes real structure on a KO-dim 3 spectral triple satisfies
$J^2 = -I$ by convention (cf.~Paper~32~\S~IV
\texttt{rem:axiom\_audit}).
\end{proof}

The two antilinear conjugations coexist on the same Hilbert spaces:\
$J_{\mathrm{GV}}$ on $\Hilb_{\nmax}$, $J_{\TT}$ on $\Hilb_{\GNS} \cong
B(\Hilb_W)$. Conflating them (e.g., by importing
\texttt{geovac.real\_structure.RealStructure} into the
modular-Hamiltonian construction) would silently produce algebraic
errors, since $J_{\TT}^2 - (-I) = 2I \ne 0$.

\begin{remark}[State-dependence vs.\ kinematic intrinsicness]
\label{rem:state_dependent}
$J_{\mathrm{GV}}$ is intrinsic to the spectral triple at construction
time:\ once $(\Acal, \Hilb, D, J, \gamma)$ is specified,
$J_{\mathrm{GV}}$ is fixed by the KO-dim convention. $J_{\TT}$ is
\emph{state-dependent}:\ each KMS state $\omega$ defines its own
$J_{\TT}^{(\omega)}$ via Tomita's theorem
on the GNS triple of $\omega$. The state-dependence is consistent
with the standard Tomita--Takesaki theorem~\cite{takesaki1970, tomita1967}.
\end{remark}

\subsection{Bit-exact period closure at the Tomita level}
\label{sec:tt_period}

\begin{theorem}[BW-$\gamma$ period closure at finite cutoff]
\label{thm:bw_gamma}
For every $\nmax \ge 1$ and every $a \in B(\Hilb_W)$, the
Tomita--Takesaki modular flow~\eqref{eq:TT_flow} satisfies the
period closure
\begin{equation}\label{eq:bw_gamma_closure}
   \sigma_{2\pi}^{\TT}(a) \;=\; a
\end{equation}
identically, with maximum periodicity residual
\begin{equation*}
   r_W^{\TT}(O)
   \;:=\; \Frob{\sigma_{2\pi}^{\TT}(O) - O}
\end{equation*}
scaling as $O(\sqrt{\dim_W} \cdot \varepsilon_{\mathrm{machine}})$ in
floating-point computation.
\end{theorem}

\begin{proof}
From~\eqref{eq:TT_flow}, $\sigma_{2\pi}^{\TT}(a) = e^{-i \cdot 2\pi
\cdot K_\alpha^W}\, a\, e^{+i \cdot 2\pi \cdot K_\alpha^W}$. By
Proposition~\ref{prop:integer_spectrum}, $K_\alpha^W$ has integer
spectrum, so $e^{\pm i \cdot 2\pi \cdot K_\alpha^W} = I$ identically.
Hence $\sigma_{2\pi}^{\TT}(a) = a$.
\end{proof}

\paragraph{Numerical verification.}
In floating-point computation, the modular operator $\Delta$ is built
on the GNS Hilbert--Schmidt space $\Hilb_{\GNS} \cong \C^{\dim_W^2}$
via the tensor-product form~\eqref{eq:Delta_matrix}, with each
$\rho^{\pm 1}$ assembled by matrix exponential of $\pm K_\alpha^W$ and
a numerical division by $Z$. The modular flow $\sigma_t^{\TT}$ is then
computed by acting $\Delta^{it}$ on each input operator $a$ written in
$\mathrm{vec}$ form, then reshaping. We have computed the maximum
periodicity residual $r_W^{\TT}(O)$ at $\nmax \in \{2, 3, 4, 5\}$;\
results are in Table~\ref{tab:bw_gamma_residuals}.

\begin{table}[ht]
\centering
\caption{Maximum BW-$\gamma$ Tomita-level periodicity residual
$\max_O r_W^{\TT}(O)$ at finite cutoff, computed via polar
decomposition on the GNS Hilbert--Schmidt space. The slightly larger
residual at $\nmax = 5$ (vs.~Table~\ref{tab:bw_alpha_residuals}) is
pure numerical conditioning from the
$\dim_W^2 \times \dim_W^2 = 4900 \times 4900$ matrix diagonalisation
of $\Delta$ with extreme eigenvalue ratios
$\lambda_i / \lambda_j$;\ no structural drift. Source data:\
\texttt{debug/data/l1\_tighten\_tomita\_results.json}.}
\label{tab:bw_gamma_residuals}
\begin{tabular}{rrrcc}
\toprule
$\nmax$ & $\dim_W$ & $\dim \Hilb_{\GNS}$ &
$\max_O r_W^{\TT}(O)$ &
verdict\\
\midrule
2 & 8  & 64   & $1.09 \times 10^{-16}$ & UNIFIED\_STRONG \\
3 & 20 & 400  & $1.07 \times 10^{-15}$ & UNIFIED\_STRONG \\
4 & 40 & 1600 & $3.03 \times 10^{-15}$ & UNIFIED\_STRONG \\
5 & 70 & 4900 & $4.79 \times 10^{-15}$ & UNIFIED\_STRONG \\
\bottomrule
\end{tabular}
\end{table}

The $J_{\TT}^2 = +I$ identity of Proposition~\ref{prop:J_TT_distinct}
is verified numerically at every $\nmax$ to machine precision
(residuals $\le 7 \times 10^{-17}$):\ this distinguishes $J_{\TT}$
from $J_{\mathrm{GV}}$ at the genuine state-dependent level, not just
at the canonical $U = I$ sanity case.

\subsection{Spectrum of $K_{\TT}$}\label{sec:K_TT_spectrum}

\begin{proposition}[Spectrum of $K_{\TT}$ on $\Hilb_{\GNS}$]
\label{prop:K_TT_spectrum}
The eigenvalues of $K_{\TT} = -\log \Delta$ on $\Hilb_{\GNS}$ are
$\{n_j - n_k : 1 \le j, k \le \dim_W\}$, where $\{n_j\}$ is the integer
spectrum of $K_\alpha^W$. In particular:
\begin{enumerate}\setlength{\itemsep}{2pt}
\item[(a)] $\mathrm{Spec}(K_{\TT}) \subset \Z$ (integer spectrum on
$\Hilb_{\GNS}$);
\item[(b)] $K_{\TT}$ has at least $\dim_W$ zero eigenvalues (from the
diagonal $j = k$ pairs);
\item[(c)] The maximum eigenvalue of $K_{\TT}$ is $2 \max_j n_j =
2(2(\nmax - 1) + 1) = 2(2\nmax - 1)$;\ for $\nmax = 2$ this is $\le 4$,
and the maximum is exactly $2$ on the wedge (since $\max_W
\mathrm{two}\_m_j = +1$ at $\nmax = 2$).
\end{enumerate}
\end{proposition}

\begin{proof}
By~\eqref{eq:Delta_matrix}, $\Delta = (\rho^{-1})^T \otimes \rho$ in
the vec representation, with $\rho = e^{-K_\alpha^W}/Z$. The
eigenvalues of $\rho$ are $\{e^{-n_j}/Z\}_{j = 1}^{\dim_W}$, so the
eigenvalues of $\Delta$ are $\{e^{n_k} \cdot e^{-n_j}\}_{j, k} =
\{e^{n_k - n_j}\}_{j, k}$. Hence $K_{\TT} = -\log \Delta$ has
eigenvalues $\{n_j - n_k\}$, which are integer differences of the
integer spectrum of $K_\alpha^W$. The diagonal $j = k$ pairs contribute
$n_j - n_j = 0$, giving at least $\dim_W$ zero eigenvalues (more if
$K_\alpha^W$ has spectral degeneracies).
\end{proof}

The integer-spectrum property
of Proposition~\ref{prop:K_TT_spectrum}(a) is the structural reason
the bit-exact closure of Theorem~\ref{thm:bw_gamma} holds:\
$e^{i \cdot 2\pi \cdot K_{\TT}} = I$ on $\Hilb_{\GNS}$ by
$e^{i \cdot 2\pi \cdot m} = 1$ for integer $m$, so the modular flow
$\Delta^{it}$ acting at $t = 2\pi$ collapses to the identity, and
$\sigma_{2\pi}^{\TT}(a) = I \cdot a \cdot I = a$ for every $a$. This is
the same mechanism as Theorem~\ref{thm:bw_alpha} but realised at the
GNS-Hilbert--Schmidt level.

\subsection{Propinquity rate cross-check}\label{sec:propinquity_consistency}

The closure of Theorem~\ref{thm:bw_gamma} is a finite-cutoff identity,
not a propinquity limit statement. To confirm consistency with
Paper~38's qualitative-rate propinquity bound~\eqref{eq:paper38_main},
we compute the ratio
\begin{equation*}
   \frac{\max_O r_W^{\TT}(O)}{\gamma_{\nmax}^{(\mathrm{L2})}}
\end{equation*}
where $\gamma_{\nmax}^{(\mathrm{L2})}$ is the L2 mass-concentration
moment of the central spectral Fej\'er kernel on $\SU(2)$
(\cite[Lemma~L2]{paper38}):\
$\gamma_2^{(\mathrm{L2})} = 2.0746$, $\gamma_3^{(\mathrm{L2})} = 1.6101$,
$\gamma_4^{(\mathrm{L2})} = 1.3223$, $\gamma_5^{(\mathrm{L2})} = 1.1302$.
The ratios are $8.7 \times 10^{-17}$, $3.0 \times 10^{-16}$,
$7.1 \times 10^{-16}$, $1.4 \times 10^{-15}$:\ several orders of
magnitude below the qualitative-rate bound at every tested $\nmax$.
The L1 closure is therefore \emph{much stronger} than the qualitative-rate
propinquity bound would predict at finite cutoff; the closure does not
require the GH limit and is structural rather than asymptotic.

% =====================================================================
\section{Unified closure}\label{sec:unified_closure}
% =====================================================================

The two constructions of Sections~\ref{sec:geometric_construction}
and~\ref{sec:tt_construction} are not the same operator on
$\Hilb_W$ versus $\Hilb_{\GNS}$ (the GNS construction lives on a
larger Hilbert space). But the modular flows they generate on the
algebra $B(\Hilb_W)$ are intimately related, and at the period
$t = 2\pi$ they collapse to the identical operator-system identity.
This section makes the relation precise.

\subsection{Flow conjugacy}\label{sec:flow_conjugacy}

\begin{theorem}[BW-$\alpha$ / BW-$\gamma$ flow conjugacy]
\label{thm:unified}
For every $a \in B(\Hilb_W)$ and every $t \in \R$,
\begin{equation}\label{eq:tt_alpha_conjugate}
   \sigma_t^{\TT}(a) \;=\; \sigma_{-t}^{\alpha}(a),
\end{equation}
where $\sigma_t^{\alpha}$ is the geometric BW-$\alpha$ flow
\eqref{eq:alpha_flow} and $\sigma_t^{\TT}$ is the Tomita--Takesaki
BW-$\gamma$ flow~\eqref{eq:TT_flow}.
\end{theorem}

\begin{proof}
By~\eqref{eq:TT_flow}, $\sigma_t^{\TT}(a) = e^{-it K_\alpha^W}\, a\,
e^{+it K_\alpha^W}$. By~\eqref{eq:alpha_flow}, $\sigma_{-t}^{\alpha}(a) =
e^{-it K_\alpha^W}\, a\, e^{+it K_\alpha^W}$. The two expressions are
literally identical.
\end{proof}

\paragraph{Consequence at the period.}
At $t = 2\pi$, equation~\eqref{eq:tt_alpha_conjugate} reads
$\sigma_{2\pi}^{\TT}(a) = \sigma_{-2\pi}^{\alpha}(a) =
\sigma_{2\pi}^{\alpha}(a)^{-1}$, where the last equality uses that the
flow $\sigma_t^{\alpha}$ is a one-parameter group with the inverse at
$t = 2\pi$ being the flow at $t = -2\pi$. But by
Theorem~\ref{thm:bw_alpha}, $\sigma_{2\pi}^{\alpha}(a) = a = a^{-1}$
trivially (this is an operator equation, not a group inverse), so
$\sigma_{2\pi}^{\TT}(a) = a$ as well, consistent with
Theorem~\ref{thm:bw_gamma}. The two closures are not just both true;
they are the \emph{same} closure under the conjugacy
\eqref{eq:tt_alpha_conjugate}.

\paragraph{Numerical confirmation at general $t$.}
At a general non-period $t$ (e.g., $t = 1$), the two flows
$\sigma_t^{\TT}$ and $\sigma_t^{\alpha}$ are distinct (the conjugacy
\eqref{eq:tt_alpha_conjugate} relates $\sigma_t^{\TT}$ to
$\sigma_{-t}^{\alpha}$, not to $\sigma_t^{\alpha}$). We have verified
this numerically:\ the average cross-flow difference
$\langle \Frob{\sigma_1^{\TT}(O) - \sigma_1^{\alpha}(O)} \rangle_O$
over the same five test operators is $\approx 0.16$ at $\nmax = 2$,
clearly non-zero. The sign-flipped comparison
$\langle \Frob{\sigma_1^{\TT}(O) - \sigma_{-1}^{\alpha}(O)} \rangle_O$
gives $\approx 4 \times 10^{-17}$, machine precision. The conjugacy
\eqref{eq:tt_alpha_conjugate} is therefore verified bit-exactly at
general $t$, not just at the period.

\subsection{The derived-Hamiltonian finding}
\label{sec:derived_hamiltonian}

The unified closure rests on the deliberate choice
$H_{\mathrm{local}} = K_\alpha^W / \beta$
(equation~\eqref{eq:H_local_def}) for the local Hamiltonian
generating the wedge KMS state. This is a construction-level choice,
not a derivation:\ we wrote down the wedge KMS state as $\rho_W =
e^{-K_\alpha^W}/Z$ rather than (say) $\rho_W = e^{-\beta D_W}/Z$ for
the wedge-restricted Camporesi--Higuchi Dirac. The motivation is
structural (we want the modular Hamiltonian's period to match the
canonical BW $\beta = 2\pi$), but it is \emph{a} choice of local
Hamiltonian, not the unique choice.

Two consequences are immediate, and together they isolate the choice
as a derived structural finding rather than a uniqueness theorem:

\paragraph{(I) Matches the Wightman BW vacuum.}
In the continuum Bisognano--Wichmann theorem, the local Hamiltonian
on the Rindler wedge is the boost generator $K_{\mathrm{boost}}$, and
the Wightman vacuum restricted to the wedge algebra is the Gibbs state
$\rho_W = e^{-2\pi K_{\mathrm{boost}}}/Z$ at inverse temperature
$\beta = 2\pi$. Our choice $H_{\mathrm{local}} = K_\alpha^W / \beta$
matches this structure exactly:\ $\beta H_{\mathrm{local}} = K_\alpha^W$
is the wedge-restricted boost-class generator (the spatial rotation
generator around the Hopf-base axis, on the Riemannian side), and the
KMS state is $\rho_W = e^{-K_\alpha^W}/Z = e^{-2\pi K_\alpha^W/(2\pi)}/Z$,
the operator-system analog of the Wightman BW vacuum. The choice is
therefore the canonical operator-system analog of the Wightman state,
not a post-hoc engineering choice.

\paragraph{(II) The framework's intrinsic Dirac is not the right
local Hamiltonian.}
If we had chosen $H_{\mathrm{local}} = D_W$ (the wedge-restricted
Camporesi--Higuchi Dirac), then $\beta H_{\mathrm{local}} = 2\pi D_W$
would have half-integer spectrum on the spinor basis (the eigenvalues
of $D_W$ are $\chi \cdot (n + 1/2)$, with $n + 1/2$ half-integer), so
$e^{i \cdot 2\pi \cdot \beta D_W}$ would \emph{not} be the identity at
$\beta = 2\pi$. The bit-exact period closure would not hold, and the
unified BW-$\alpha$ / BW-$\gamma$ structure would split.

This is a non-trivial structural finding. The framework's intrinsic
Dirac $\DCH$ is the canonical spectral-action object on the
Camporesi--Higuchi triple, and one might naively expect it to play the
role of the local Hamiltonian for any natural KMS state on the
truncated operator system. The above shows that this is not the case
at $\beta = 2\pi$:\ the right local Hamiltonian for the BW vacuum on
the wedge is the boost-class generator $K_\alpha^W$, not the
wedge-restricted Dirac $D_W$. The two play complementary roles:\
$\DCH$ is the spectral-action object (the Connes--Chamseddine input
for the heat-kernel expansion); $K_\alpha^W$ is the
wedge-modular-Hamiltonian object (the Tomita--Takesaki input for the
modular-flow construction). On the round $\sthree$ they happen to
agree up to an overall scale (both are diagonal in the spinor basis
with simple spectral structure), but only one of them has the integer
spectrum required for the bit-exact period closure at $\beta = 2\pi$.

\begin{remark}[A speculative reading]\label{rem:speculation}
We note, but do not pursue here, the possibility that the
distinction between the spectral-action Dirac $\DCH$ and the
modular-Hamiltonian generator $K_\alpha^W$ is structurally significant
for the framework's interpretation. In the spectral-action paradigm
(Chamseddine--Connes~\cite{chamseddine_connes2010}), the Dirac is the
universal object from which all physics is read; in the
thermal-time paradigm (Connes--Rovelli~\cite{connes_rovelli1994}), the
modular Hamiltonian of an algebraic state is the physical-time generator.
On the round $\sthree$ Camporesi--Higuchi triple at the wedge KMS
state $\rho_W = e^{-K_\alpha^W}/Z$, these are not the same operator.
The structural status of this distinction --- whether it is a
peculiarity of the round $\sthree$ truncation or a deeper feature of
the framework's spectral content --- is open. We flag it as one of
the open questions in Section~\ref{sec:conclusion}.
\end{remark}

\subsection{Cross-witness collapse at the algebra-action level}
\label{sec:cross_witness_algebra}

The four-witness collapse mentioned in
Section~\ref{sec:four_witness_unit} acquires a sharper form at the
algebra-action level. With $H_{\mathrm{local}} = K_\alpha^W / \beta$,
the density matrix is
\begin{equation}\label{eq:rho_W_beta_indep}
   \rho_W \;=\; \frac{e^{-\beta H_{\mathrm{local}}}}{Z}
            \;=\; \frac{e^{-K_\alpha^W}}{Z}
\end{equation}
\emph{independent of $\beta$}. The Tomita modular operator $\Delta$
constructed from $\rho_W$ via~\eqref{eq:Delta_matrix} and the modular
Hamiltonian $K_{\TT} = -\log \Delta$ are therefore also
$\beta$-independent. The four-witness theorem at the operator-system
level reduces to the assertion that all four physical witnesses
(BW, HH, Sew, Unruh) produce the \emph{same} modular operator
$\Delta$ and the same modular Hamiltonian $K_{\TT}$, with the
witness-specific $\kappa_g$ values dropping out of the algebra-action
construction.

We have verified this numerically at $\nmax \in \{2, 3, 4, 5\}$ across
the six witness instantiations (BW canonical at $\kappa_g = 1$;
Hartle--Hawking at $M = 1$, $\kappa_g = 1/4$; Hartle--Hawking at
$M = 2$, $\kappa_g = 1/8$; Sewell at $M = 1$, $\kappa_g = 1/4$;
Unruh at $a = 1$, $\kappa_g = 1$; Unruh at $a = 2$, $\kappa_g = 2$).
The cross-witness consistency residual is exactly $0$ at every tested
$\nmax$ between any pair of witnesses, confirming the
$\beta$-independence of the density matrix at the construction level.

% =====================================================================
\section{Four-witness collapse and a single construction}
\label{sec:four_witness_collapse}
% =====================================================================

The numerical content of the four-witness collapse is summarised in
Table~\ref{tab:cross_witness} and the structural content is the
following corollary.

\begin{table}[ht]
\centering
\caption{Cross-witness consistency at the BW-$\alpha$ level
(left columns) and at the BW-$\gamma$ Tomita level (right columns).
Cross-consistency residual is the difference between the period
residual of each witness and the BW canonical period residual at the
same $\nmax$:\ exactly $0$ across all six witnesses
at every tested $\nmax$. Source data:\
\texttt{debug/data/l1\_modular\_hamiltonian\_results.json},
\texttt{debug/data/l1\_tighten\_tomita\_results.json}.}
\label{tab:cross_witness}
\renewcommand{\arraystretch}{1.1}
\setlength{\tabcolsep}{4pt}
\begin{tabular}{rcccccc}
\toprule
$\nmax$ & $\BW$ & $\HH_{M=1}$ & $\HH_{M=2}$ & $\Sew_{M=1}$ & $\Unruh_{a=1}$ & $\Unruh_{a=2}$ \\
& ($\kappa_g\!=\!1$) & ($\kappa_g\!=\!1/4$) & ($\kappa_g\!=\!1/8$) & ($\kappa_g\!=\!1/4$) & ($\kappa_g\!=\!1$) & ($\kappa_g\!=\!2$) \\
\midrule
\multicolumn{7}{l}{\emph{BW-$\alpha$ residual (Table~\ref{tab:bw_alpha_residuals}):}} \\
2 & 1.80e-16 & 1.80e-16 & 1.80e-16 & 1.80e-16 & 1.80e-16 & 1.80e-16 \\
3 & 4.76e-16 & 4.76e-16 & 4.76e-16 & 4.76e-16 & 4.76e-16 & 4.76e-16 \\
4 & 9.33e-16 & 9.33e-16 & 9.33e-16 & 9.33e-16 & 9.33e-16 & 9.33e-16 \\
5 & 1.59e-15 & 1.59e-15 & 1.59e-15 & 1.59e-15 & 1.59e-15 & 1.59e-15 \\
\midrule
\multicolumn{7}{l}{\emph{BW-$\gamma$ Tomita residual (Table~\ref{tab:bw_gamma_residuals}):}} \\
2 & 1.09e-16 & 1.09e-16 & 1.09e-16 & 1.09e-16 & 1.09e-16 & 1.09e-16 \\
3 & 1.07e-15 & 1.07e-15 & 1.07e-15 & 1.07e-15 & 1.07e-15 & 1.07e-15 \\
4 & 3.03e-15 & 3.03e-15 & 3.03e-15 & 3.03e-15 & 3.03e-15 & 3.03e-15 \\
5 & 4.79e-15 & 4.79e-15 & 4.79e-15 & 4.79e-15 & 4.79e-15 & 4.79e-15 \\
\bottomrule
\end{tabular}
\end{table}

\begin{corollary}[Single operator-system construction for the four
witnesses]\label{cor:single_construction}
Under the unit normalisation U-1 and the wedge KMS state choice
$\rho_W = e^{-K_\alpha^W}/Z$, the operator-system-level realisations of
the Hartle--Hawking, Sewell, Bisognano--Wichmann, and Unruh witnesses
on the truncated Camporesi--Higuchi triple $\Tcal_{\nmax}$ collapse to
a single construction. The modular operator $\Delta$ and the modular
Hamiltonian $K_{\TT}$ are bit-identical across the six witness
instantiations. The witness-specific physics enters only through the
parameterisation of $\kappa_g$, which is invisible at the
algebra-action level and at the period closure $\sigma_{2\pi}(O) = O$.
\end{corollary}

\begin{proof}
The density matrix $\rho_W = e^{-K_\alpha^W}/Z$ is
$\beta$-independent (Section~\ref{sec:cross_witness_algebra}). The
Tomita modular operator $\Delta$ on the GNS Hilbert--Schmidt space is
constructed entirely from $\rho_W$ via~\eqref{eq:Delta_matrix},
and the modular Hamiltonian $K_{\TT} = -\log \Delta$ depends only on
$\Delta$. Both are therefore $\beta$-independent.
$\beta$-independence at the algebra-action level immediately implies
identical period closure at $t = 2\pi$ for all six witness
instantiations.
\end{proof}

\paragraph{Structural reading.}
At the operator-system level on the truncated Camporesi--Higuchi
triple, the four-witness Wick-rotation theorem reduces to one
construction. The witness-specific physical content (mass of the black
hole, proper acceleration of the Rindler observer) parameterises the
correspondence to the continuum physical observable but does not
modify the underlying modular-Hamiltonian construction. The same
$2\pi$ that appears as the Hopf-base measure factor $\Vol(S^2)/\pi^2$
in the Paper~38 propinquity rate constant
\eqref{eq:paper38_main} also drives the modular period closure of
Theorems~\ref{thm:bw_alpha}--\ref{thm:bw_gamma}.

This is the operator-system-level realisation of the M1 sub-mechanism
of the master Mellin engine
(Paper~32~\S~VIII case-exhaustion theorem of the present author's
framework). The metric-functional-level reading
(\cite[Rem.~rem:bisognano\_wichmann\_reading]{paper32}, Sprint TD
Track 4, Sprint Unruh-pendant) identifies the framework's $2\pi$
on a temporal $S^1_\tau$ with the modular-flow / boost-orbit period
of the four witnesses via the published Wick-rotation chain. The
present paper lifts that identification to the operator-system level
inside the spectral triple, eliminating the dependence on the
Wick-rotation prescription for the period-closure identity itself:\
the $2\pi$ in the modular period is now a framework-internal output
of the spectral truncation, not an externally imposed Wick-rotation
matching.

\begin{remark}[Pythagorean orthogonality underlies the six-witness
collapse]\label{rem:pythagorean_underlies_collapse}
A post-construction PSLQ probe on the Lorentzian extension
(\cite[\S~10.2,
Cor.~cor:pythagorean\_orthogonality]{paper43};\
\texttt{debug/h\_local\_residual\_pslq\_memo.md},
\texttt{debug/six\_witness\_hs\_orthogonality\_memo.md})
identified a structural inner-product identity that underlies the
six-witness collapse of
Corollary~\ref{cor:single_construction}. For every witness
$w \in \{\BW, \HH_{M=1}, \HH_{M=2}, \Sew_{M=1}, \Unruh_{a=1},
\Unruh_{a=2}\}$, the BW-aligned local Hamiltonian
$H_{\mathrm{local}}^{(w)} := K_\alpha^W / \beta^{(w)}$ and the
wedge-restricted Camporesi--Higuchi Dirac $D_W^{\mathrm{GV}}$ are
Hilbert--Schmidt orthogonal:
$$
   \big\langle H_{\mathrm{local}}^{(w)},\, D_W^{\mathrm{GV}} \big\rangle_{\mathrm{HS}}
   \;:=\; \Tr\!\big( (H_{\mathrm{local}}^{(w)})^\dagger\, D_W^{\mathrm{GV}} \big)
   \;=\; 0,
$$
verified bit-exactly at every cell of the 18-cell panel (six
witnesses $\times$ $\nmax \in \{2, 3\}$;\ extended to $\nmax \in
\{1, \ldots, 5\}$ on the BW canonical row;\ max $|\langle H, D
\rangle_{\mathrm{HS}}| = 8.9\times 10^{-16}$, machine precision).
The witness-independence of the inner product is a
$\kappa_g$-linearity corollary:\ $H_{\mathrm{local}}^{(w)}$ scales
linearly with $\kappa_g^{(w)}$, while $D_W^{\mathrm{GV}}$ is
witness-independent.

The mechanism is structural rather than computational:\ in the
unfolded wedge spinor basis, $H_{\mathrm{local}}^{(w)}$ is real and
diagonal (eigenvalues $\mathrm{two\_m}_j > 0$, the rotation
generator around the Hopf axis), while $D_W^{\mathrm{GV}}$ is
off-diagonal (it intertwines $\pm m_j$ chirality partners via the
wedge basis change). The two operators occupy mutually-orthogonal
subspaces of $\mathcal{B}(\Hilb_W)$, so the inner product vanishes;\
the diagonal/off-diagonal subspace decomposition is sketched at
the basis level and supports a Pythagorean form for the residual
$$
   \big\|H_{\mathrm{local}}^{(w)} - D_W^{\mathrm{GV}}\big\|_F^2
   \;=\; \frac{(\kappa_g^{(w)})^2 \cdot S(\nmax)}{4\pi^2}
       \;+\; D(\nmax),
$$
with $S, D$ pure rationals in $\nmax$ (Paper~43~\S~10.2). A formal
proof of the subspace decomposition at general $\nmax$ on the
truncated spectral triple is a named follow-on; the empirical
verification across the 18-cell panel is strong evidence the
orthogonality is structural rather than coincidental.

\emph{Reading.}\ Corollary~\ref{cor:single_construction} states
that the six witnesses share a single modular operator and a
single modular Hamiltonian via $\beta$-independence of $\rho_W$.
The present remark sharpens that reading at the inner-product
level:\ the period closure $\sigma_{2\pi}(O) = O$ for every $O$
in the operator system is underlaid by a structural Pythagorean
identity that controls how $H_{\mathrm{local}}^{(w)}$ and
$D_W^{\mathrm{GV}}$ decompose orthogonally in HS norm, with the
witness-specific physical content carried by the
$\kappa_g$-dependence of the M1 piece and not by any rotation
between the two subspaces.
\end{remark}

% =====================================================================
\section{Related work and scope}\label{sec:related_work}
% =====================================================================

\paragraph{Connes--Rovelli thermal time hypothesis.}
The closest published analog of the present construction is the
Connes--Rovelli thermal time hypothesis~\cite{connes_rovelli1994},
which identifies the Tomita modular flow of a general-covariant von
Neumann algebra with the physical-time evolution. The present paper
specialises the Connes--Rovelli construction to a finite-dimensional
Type I$_{\dim_W}$ factor obtained as the wedge-restricted
Connes--van~Suijlekom~\cite{connes_vs2021} truncation of the round
$\sthree$ Camporesi--Higuchi spectral triple. The Connes--Rovelli
construction is abstract;\ the present construction is concrete in
the sense that $\Delta$, $K_{\TT}$, and the period closure are all
explicit finite-dimensional matrix operations, computable in
floating-point arithmetic at any tested $\nmax$.

\paragraph{Bisognano--Wichmann tradition and continuum modular
Hamiltonians.}
The original Bisognano--Wichmann
papers~\cite{bisognano_wichmann1975, bisognano_wichmann1976} establish
the modular-flow-as-boost identification for the Rindler wedge of
Minkowski spacetime, with the analytic period $\sigma_{i \cdot 2\pi} =
\mathrm{id}$ as a corollary of the Tomita--Takesaki KMS
condition~\cite{takesaki1970, tomita1967} at $\beta = 2\pi$ in
rapidity units. The continuum entanglement-Hamiltonian construction
for a half-space in free CFT was developed by Casini and
Huerta~\cite{casini_huerta2009}, and extended to spherical regions of
$S^d$ in CFT by Casini--Huerta--Myers~\cite{casini_huerta_myers2011}.
The present construction is the operator-system-level analog of the
Casini--Huerta--Myers spherical-region modular Hamiltonian on the
truncated Camporesi--Higuchi triple over the round $\sthree$, with
the integer spectrum of $K_\alpha^W$ playing the role of the
canonical $2\pi$-period of the continuum modular flow.

\paragraph{Lattice realisations.}
Zhu--Casini~\emph{et~al.}~\cite{zhu_casini2020} verified the
Bisognano--Wichmann form of the modular Hamiltonian on critical
quantum spin chains, finding the continuum BW prediction to hold to
good approximation on finite lattices. The present paper differs in
two respects:\ (i)~the underlying object is a continuum spectral
triple on a non-flat compact $\sthree$ background, not a lattice
discretisation of physical space; (ii)~the finite-dimensional
structure comes from the spectral truncation of $\DCH$ via the
Connes--van~Suijlekom framework, not from a real-space discretisation.
The structural mechanism (integer-spectrum generator forcing the
$\sigma_{2\pi} = \mathrm{id}$ period closure) is morally analogous to
the lattice realisations of Zhu et al., but the construction lives on
the operator-system side of the metric-spectral-triple framework
rather than on the lattice-Hamiltonian side of condensed-matter
physics.

\paragraph{Algebraic QFT on curved backgrounds.}
The standard algebraic-QFT references for modular structures on
Riemannian backgrounds are
Strohmaier~\cite{strohmaier2006} (Reeh--Schlieder on stationary
spacetimes) and Verch~\cite{verch2001} (nuclearity-modular on compact
Riemannian backgrounds). These address the modular-theoretic content
of axiomatic QFT on curved backgrounds. The present construction is
spectral-triple-internal:\ it does not invoke a Wightman field algebra
or a nuclearity argument, but rather works directly with the
$\Op_{\nmax} \subset B(\Hilb_{\nmax})$ truncation of the
$C^\infty(\sthree) \subset B(L^2(\sthree, \Sigma))$ continuum algebra
and the associated Hilbert--Schmidt GNS construction.

\paragraph{Lorentzian signature and Krein-space spectral triples.}
The proper Lorentzian extension of the present construction would
require a Krein-space spectral triple at signature $(3, 1)$, for which
the canonical reference is
Bizi--Brouder--Besnard~2018~\cite{bizi_brouder_besnard2018}. The
Bizi--Brouder--Besnard classification organises spectral triples by
pairs $(m, n) \in (\Z/8)^2$ where $m + n$ is the metric dimension and
the underlying Connes axiom signs determine the Krein-space structure.
The Riemannian Camporesi--Higuchi triple sits at $(3, 0)$;\ the
Lorentzian extension would sit at $(3, 1)$. The
Bizi--Brouder--Besnard Table~2 signs at $(3, 1)$ flip the
$\{J, \gamma_5\}$ anticommutation relation, which leads to the
present author's Sprint L0 prediction that the M3 sub-mechanism of
the master Mellin engine trivialises at $(3, 1)$ on chirality-symmetric
spectrum (\texttt{debug/lorentzian\_l0\_audit\_memo.md} \S 3).
Constructing the Krein-space lift, replicating the present paper's
modular construction at $(3, 1)$, and verifying the M3 trivialisation
prediction is the named Sprint L2 follow-up of the present author's
framework documentation:\ a multi-month effort dominated by the
absence of a published Lorentzian analog of the Latr\'emoli\`ere
propinquity (a NCG-mathematics gap independent of any specific physical
application).

\paragraph{Twisted-spectral-triple emergence.}
The recent work of
Nieuviarts~\cite{nieuviarts2025a, nieuviarts2024, nieuviarts2025b_proceedings}
proposes a twisted-spectral-triple morphism that constructs a
pseudo-Riemannian spectral triple from a Riemannian almost-commutative
twisted spectral triple, presented as an algebraic alternative to
Wick rotation. The construction does not apply to the present
$\sthree$ case because of the structural odd-dimensional restriction
(\S~\ref{sec:nieuviarts_footnote}). The Lorentzian-extension route for
the present paper therefore runs through the direct
Bizi--Brouder--Besnard~2018 Krein-space lift, not through a twist
morphism.

\paragraph{Other NCG-side gaps the present paper does not close.}
A genuine Lorentzian propinquity (in the sense of an analog of the
Latr\'emoli\`ere quantum Gromov--Hausdorff propinquity for
Krein-space metric spectral triples) is open as of May 2026.
Toyota~\cite{toyota2023}, Farsi--Latr\'emoli\`ere~\cite{farsi_latremoliere2024,
farsi_latremoliere2025}, Hekkelman--McDonald--van~Suijlekom~\cite{ucp_maps_2024},
and Hekkelman--McDonald~\cite{hekkelman_mcdonald2024b} all extend
the metric-spectral-triple propinquity machinery in various directions
(discrete groups of polynomial growth, collapse with abelian factors,
analytic paths of Riemannian metrics, noncommutative integral on
truncations, Berezin--Toeplitz on $S^2$), but none addresses
Krein-space signature change. Constructing the Lorentzian propinquity
is original NCG-mathematics work, not algorithmic.

% =====================================================================
\section{Conclusion and open questions}\label{sec:conclusion}
% =====================================================================

We have established two structurally independent but operator-action
conjugate constructions of the modular Hamiltonian on the truncated
Camporesi--Higuchi spectral triple $\Tcal_{\nmax}$, both verifying
the four-witness Wick-rotation period closure
$\sigma_{2\pi}(O) = O$ \emph{bit-exactly} at every tested
$\nmax \in \{2, 3, 4, 5\}$ and all four physical witnesses
(Bisognano--Wichmann, Hartle--Hawking, Sewell, Unruh) under the
canonical hemispheric wedge $P_W$ aligned with the Hopf-base axis,
the BW local Hamiltonian $H_{\mathrm{local}} = K_\alpha^W / \beta$
on the wedge sub-algebra, and the unit normalisation $\kappa_g = 1$
in natural rapidity units. The geometric BW-$\alpha$ realisation
uses the integer-spectrum boost-class generator $K_\alpha^W =
J_{\mathrm{polar}}$ (Theorem~\ref{thm:bw_alpha}); the
Tomita--Takesaki BW-$\gamma$ realisation uses the polar
decomposition $S = J_{\TT} \Delta^{1/2}$ on the GNS Hilbert--Schmidt
space (Theorem~\ref{thm:bw_gamma}). The two flows are operator-action
conjugates, $\sigma_t^{\TT}(O) = \sigma_{-t}^{\alpha}(O)$
(Theorem~\ref{thm:unified}), and the four physical witnesses collapse
to a single construction at the algebra-action level
(Corollary~\ref{cor:single_construction}).

This lifts the four-witness Wick-rotation theorem from
``structural correspondence'' (the published Wick-rotation chain at
the metric-functional level) to \emph{literal identification at the
operator-system level} on the Riemannian truncated triple, at every
finite cutoff. The structural reason is the integer spectrum of the
BW local Hamiltonian, which forces $e^{i \cdot 2\pi \cdot K_\alpha^W}
= I$ on the underlying Hilbert space and hence the bit-exact closure
of the modular flow at $t = 2\pi$. Unlike Paper~38's qualitative-rate
propinquity convergence~\eqref{eq:paper38_main}, the present closure
does not require a Gromov--Hausdorff limit and is not asymptotic in
$\nmax$;\ it is a finite-cutoff exact identity.

\paragraph{Open questions named by the closure.}

\textbf{(O1) Lorentzian extension at signature $(3, 1)$.}
The natural next step is to replicate the present operator-system
construction on a Bizi--Brouder--Besnard Krein-space spectral triple
at $(3, 1)$, verify the Sprint L0 prediction that M3 trivialises but
M1 stays under the BBB Table~2 sign flip, and produce the genuine
Lorentzian KMS-Tomita-Takesaki theorem inside the framework. The
multi-month obstruction is the absence of a Lorentzian propinquity in
the published literature.

\emph{Progress update (Sprints L2-B/C/D, 2026-05-16).}
The Hilbert-space and Lorentzian-Dirac levels of the $(3, 1)$
extension have landed:\ Sprint L2-B
(\texttt{geovac/krein\_space\_construction.py}) constructs the
Krein space $\mathcal{K}_{\nmax, N_t} = \mathcal{H}_{\rm GV}^{\nmax}
\otimes L^2(\mathbb{R}_t)_{\rm cutoff}$ with fundamental symmetry
$J = \gamma^0$ in the Peskin--Schroeder chiral basis, verifying all
Krein axioms ($J^2 = +I$, $J^* J = I$, $J^* = J$,
$\mathcal{K} = \mathcal{K}^+ \oplus \mathcal{K}^-$) bit-exactly at
$(\nmax, N_t) \in \{1, 2, 3\} \times \{1, 11, 21\}$ and passing the
load-bearing Riemannian-limit recovery at $N_t = 1$ bit-identically;
Sprint L2-C (\texttt{geovac/lorentzian\_dirac.py}) constructs
$D_L = i \cdot [\gamma^0 \otimes \partial_t + D_{\rm GV} \otimes I_{N_t}]$
via van~den~Dungen Proposition~4.1 with $\partial_t$ centered FD plus
Dirichlet zero BC (structurally forced for anti-Hermiticity and
Krein-self-adjointness), verifying Krein-self-adjointness
$D_L^\times = D_L$ and Riemannian-limit recovery $D_L|_{N_t = 1} =
i D_{\rm GV}$ bit-exactly across the same panel.  Sprint L2-D
(\texttt{geovac/connes\_axiom\_audit\_31.py}) verifies four of the
six BBB-predicted axioms at $(m, n) = (4, 6)$ bit-exactly
($J_L^2 = +I$, $\{J_L, \gamma^5\} = 0$, $\{J_L, \eta\} = 0$,
$J_L D_L = +D_L J_L$); the fifth axiom $\{\gamma^5, D_L\} = 0$ is a
\emph{structural finding}, failing on the truthful Camporesi--Higuchi
Dirac (which is chirality-diagonal in the chiral basis and therefore
commutes with $\gamma^5$) and requiring an R1+R2 split documented in
Paper~32~\S~IV / \S~VIII.D~\cite{paper32}.  This is structurally
analogous to the Riemannian-side R3.5/present-paper trade between
truthful CH (correct for axioms requiring $JD = +DJ$) and offdiag CH
(correct for SDP-bounding Connes-distance constructions).  The
construction does \emph{not} re-open WH1 PROVEN
(Paper~38~\cite{paper38}); the Krein construction is an extension on
top of the Riemannian foundation, not a re-test of it.  Full
operator-system-level closure (the analog of the present paper's
literal identification of the four-witness theorem at Lorentzian
signature) is the named Sprint L2-E target, still open at the
multi-month scale due to the absence of a published Lorentzian
propinquity.  O1 status remains OPEN; the Hilbert-space and
Dirac-operator levels are now closed at finite cutoff.

\textbf{(O2) Non-tracial wedge states.}
The unified closure rests on the tracial-Gibbs structure
$\rho_W = e^{-K_\alpha^W}/Z$ on a finite-dimensional Type I$_{\dim_W}$
factor. For non-tracial wedge states (e.g., excited states above the
BW vacuum, mixed states with non-Gibbs structure), the polar
decomposition $S = J_{\TT} \Delta^{1/2}$ still goes through, but the
integer-spectrum property of $K_{\TT}$ is not guaranteed, and the
bit-exact closure may degrade to a soft propinquity-rate-controlled
closure of the type Paper~38 supplies. The structural status of the
non-tracial case is open.

\textbf{(O3) Spectral-action Dirac vs.\ modular-Hamiltonian generator.}
Section~\ref{sec:derived_hamiltonian} flagged that the framework's
intrinsic Camporesi--Higuchi Dirac $\DCH$ is \emph{not} the right
local Hamiltonian for the BW vacuum at $\beta = 2\pi$:\ the BW vacuum
selects $K_\alpha^W$, not $\DCH$, as its modular generator. On the
round $\sthree$ this distinction is mostly a matter of choosing the
right generator from a family of admissible candidates, but in
general the spectral-action $\DCH$ and the modular-Hamiltonian
$K_\alpha^W$ may not even commute. The structural status of this
distinction in extensions to other compact-Lie-group spectral triples
(per Paper~40~\cite{paper40_unified}) is open and an interesting
direction for follow-up.

\emph{Refinement (Sprint L2-E, 2026-05-16; \S~\ref{sec:lorentzian_closure} below).}\ \ %
The O3 question of signature-dependence is now \emph{partially}
closed:\ the $H_{\mathrm{local}} \ne D_W$ distinction is
\textbf{signature-INDEPENDENT at the Riemannian limit} and refined
upward at $N_t > 1$ by temporal-derivative content.  At $N_t = 1$
on the Lorentzian Krein space, $\|H_{\mathrm{local}} - D_L^W\|_F$
equals the Riemannian-side residual $\|H_{\mathrm{local}} - D_{\mathrm{GV}}^W\|_F$
\emph{bit-exact} at every tested $n_{\max} \in \{1, 2, 3\}$.  The
distinction is therefore NOT a peculiarity of the round-$\sthree$
Riemannian truncation but a deeper structural feature of the
framework's modular content.  At $N_t > 1$ the Lorentzian Dirac
$D_L = i(\gamma^0 \otimes \partial_t + D_{\mathrm{GV}} \otimes I)$
has temporal-derivative content that $K_\alpha^W / \beta$ cannot
capture; the gap between the spectral-action and modular-Hamiltonian
generators widens with temporal cutoff.  The \emph{deeper} question
(why is the spectral-action Dirac not the modular generator on the
BW vacuum, structurally?) remains open and is the residual O3.

\emph{Further refinement (Pythagorean orthogonality,
2026-05-17; Remark~\ref{rem:pythagorean_underlies_collapse} above,
Paper~43~\S 10.2~\cite{paper43}).}\ \ %
A post-L2-E PSLQ probe has identified the structural mechanism
behind the signature-independent residual:\ $H_{\mathrm{local}}$
and $D_W^L$ are Hilbert--Schmidt orthogonal,
$\langle H_{\mathrm{local}}, D_W^L \rangle_{\mathrm{HS}} = 0$
bit-exact across the entire 18-cell six-witness panel, and the
squared residual decomposes as a Pythagorean sum
$\|H_{\mathrm{local}} - D_W^L\|_F^2 = \kappa_g^2 S(n_{\max}) /
(4\pi^2) + D(n_{\max})$ with $S, D$ pure rationals in $n_{\max}$.
The sole transcendental ($1/\pi^2$) places the residual entirely
in the Paper~32~\S~VIII master Mellin engine M1 sector. The
mechanism is the diagonal/off-diagonal subspace decomposition of
$\mathcal{B}(\Hilb_W)$:\ $H_{\mathrm{local}}$ is diagonal in the
unfolded wedge spinor basis, $D_W^L$ is off-diagonal. The
$S^3$-specificity of this construction --- and its failure on the
Bargmann--Segal $S^5$ Hardy sector, which lacks a spinor sector
and a half-integer wedge --- is the fourth layer of the Paper~24
\S~5.3~\cite{paper24} Coulomb/HO asymmetry (alongside
spectrum-computing $L_0$, calibration $\pi$, and Wilson matter
coupling). The structural-explanation question (\emph{why} is
$D_W$ the wrong modular generator?) is now sharpened to a
subspace-orthogonality statement;\ a formal proof of the
diagonal/off-diagonal decomposition at general $n_{\max}$ is a
named follow-on.

\textbf{(O4) Higher-rank compact non-abelian extensions.}
Paper~40~\cite{paper40_unified} extends the Paper~38 propinquity
convergence theorem to all compact connected Lie groups with
bi-invariant metric. The natural extension of the present paper would
verify that the geometric BW-$\alpha$ generator analog on higher-rank
groups (e.g., $K_\alpha^W = J_z$ with $J_z$ the Cartan generator
along the wedge-defining direction) still has integer spectrum on the
spinor harmonic basis, preserving the bit-exact period closure across
the universal $4/\pi$-rate class. The standard structure of weights
in compact Lie groups suggests this extension is straightforward;
verifying it explicitly is a near-term sprint follow-up.

\textbf{(O5) Cross-manifold modular structures.}
The Paper~24~\cite{paper24} cross-manifold blocker
$\Tcal_{\sthree} \otimes \Tcal_{\mathrm{Hardy}(S^5)}$ remains
unaddressed by the present paper. A cross-manifold modular construction
(modular Hamiltonian on a tensor product of a round $\sthree$
Camporesi--Higuchi triple and a holomorphic-sector $S^5$
Bargmann--Segal triple) would require an analog of the Latr\'emoli\`ere
tensor-product propinquity of Paper~39~\cite{paper39} extended to
mixed Riemannian / Hardy-sector spectral triples. This is also open.

\paragraph{Final remark on the M1 master Mellin engine identification.}
The $2\pi$ in the modular period of the present construction is the
same $2\pi$ that drives the Paper~38 propinquity rate
$4/\pi = \Vol(S^2)/\pi^2$ via the M1 Hopf-base measure mechanism
(Paper~32~\S~VIII case-exhaustion theorem). The present paper
establishes this identification at the operator-system level, not
just at the metric-functional level:\ the modular period
$\sigma_{2\pi}(O) = O$ is a framework-internal output of the spectral
truncation, with the integer-spectrum generator $K_\alpha^W$ acting
as the mechanism through which the M1 transcendental signature
$\Vol(S^1) = 2\pi$ is realised at the algebra-action level. The
framework's M1 mechanism on the round $\sthree$ Camporesi--Higuchi
triple is now operator-system-level identified with the four-witness
Wick-rotation period structure, in both the geometric and the
Tomita--Takesaki readings, at every finite cutoff.

% =====================================================================
\section{Lorentzian closure at finite cutoff (Sprint L2)}
\label{sec:lorentzian_closure}
% =====================================================================

The body of this paper closes the four-witness Wick-rotation theorem
at the operator-system level on the Riemannian truncated triple
$\Tcal_{\nmax}$.  Open question (O1) of \S\ref{sec:conclusion} named
the Lorentzian extension to signature $(3, 1)$ as the next step.
Sprint L2 (May~2026) closes that extension at finite cutoff via
five sub-sprints (L2-A scoping, L2-B Krein space, L2-C Lorentzian
Dirac, L2-D Connes axiom audit, L2-E modular Hamiltonian).  This
section documents the closure and the substantive structural finding
it produced; the math.OA-style standalone writeup is outlined in
Paper~43 (\texttt{papers/standalone/paper\_43\_lorentzian\_extension\_outline.md},
companion).

\subsection{The five-sub-sprint architecture}
\label{sec:lorentzian_closure_architecture}

The Sprint L2 architecture mirrors the Riemannian construction of
\S\S\ref{sec:setup}--\ref{sec:four_witness_collapse} verbatim with a
temporal slot added.  At signature $(3, 1)$ West-coast in the
Peskin--Schroeder chiral basis (Sprint L2-B convention lock-in):

\begin{itemize}\setlength{\itemsep}{2pt}
\item \textbf{Krein space (L2-B).}\ \ %
$\mathcal{K}_{\nmax, N_t} := \mathcal{H}_{\rm GV}^{\nmax} \otimes
\mathbb{C}^{N_t}$, with $N_t$ a bounded uniform temporal grid on
$[-T_{\max}, T_{\max}]$ (NOT compactified to $S^1_\beta$, which would
re-introduce closed timelike curves per Geroch's theorem).  Fundamental
symmetry $J = \gamma^0 \otimes I_{N_t}$ with $\gamma^0$ the chirality-swap
in the chiral basis.  All Krein axioms ($J^2 = +I$, $J^* J = I$,
$J^* = J$, $\mathcal{K} = \mathcal{K}^+ \oplus \mathcal{K}^-$
orthogonally with $\dim \mathcal{K}^\pm = \dim \mathcal{K}/2$) are
bit-exact (Frobenius residual $= 0.0$ in float64) at every
$(\nmax, N_t) \in \{1, 2, 3\} \times \{1, 11, 21\}$.

\item \textbf{Lorentzian Dirac (L2-C).}\ \ %
$D_L := i \cdot \bigl[\gamma^0 \otimes \partial_t + D_{\rm GV} \otimes I_{N_t}\bigr]$
per van~den~Dungen Proposition~4.1 (arXiv:1505.01939) with
$\partial_t$ centered finite-difference plus Dirichlet zero BC
(structurally forced for anti-Hermiticity).  Krein-self-adjointness
$D_L^\times = D_L$ is bit-exact across the panel; the sign $i^t = +i$
for $t = 1$ is derived from the Krein-self-adjointness requirement,
not chosen.

\item \textbf{Connes axiom audit (L2-D).}\ \ %
At Bizi--Brouder--Besnard~\cite{bizi_brouder_besnard2018} signature
$(m, n) = (4, 6)$ (the BBB Table~3 translation of Lorentzian $(s, t) = (3, 1)$
West-coast), four BBB-predicted-sign axioms hold bit-exactly:
$J_L^2 = +I$ ($\varepsilon = +1$); $\{J_L, \gamma^5\} = 0$
($\varepsilon'' = -1$); $\{J_L, \eta\} = 0$ ($\varepsilon \kappa = -1$);
and the BBB universal $J_L D_L = +D_L J_L$.  See Paper~32~\S~VIII.D~\cite{paper32}
for the full audit.

\item \textbf{Modular Hamiltonian (L2-E, this closure).}\ \ %
Built on the wedge $W_L := P_W^{\mathrm{spatial}} \otimes P_{t \geq 0}$,
with $P_W^{\mathrm{spatial}}$ the Definition~\ref{def:wedge}
hemispheric wedge on $\mathcal{H}_{\rm GV}$ and $P_{t \geq 0}$ the
$t \geq 0$ half-line projector.  Wedge KMS state $\rho_W^L = e^{-K_L^{\alpha, W}}/Z$
under the BW choice $H_{\mathrm{local}} := K_L^{\alpha, W}/\beta$
with $\beta = 2\pi$.  The construction mirrors
\S\S\ref{sec:geometric_construction}--\ref{sec:tt_construction}
of the Riemannian closure, with the temporal slot adding additional
content at $N_t > 1$.
\end{itemize}

The architecture is structurally additive on top of the Riemannian
closure of this paper:\ at $N_t = 1$, every Sprint L2 construction
reduces \emph{bit-identically} to the corresponding Riemannian-side
object of \S\S\ref{sec:setup}--\ref{sec:four_witness_collapse}.

\subsection{The Krein-level four-witness theorem at finite cutoff}
\label{sec:krein_four_witness_theorem}

\begin{theorem}[Krein-level four-witness Wick-rotation theorem at finite cutoff]
\label{thm:krein_four_witness}
On the Lorentzian Krein space $\mathcal{K}_{\nmax, N_t}$ of
\S\ref{sec:lorentzian_closure_architecture}, with the wedge
$W_L = P_W^{\mathrm{spatial}} \otimes P_{t \geq 0}$, the wedge
KMS state $\rho_W^L = e^{-K_L^{\alpha, W}}/Z$ under the BW choice
$H_{\mathrm{local}} = K_L^{\alpha, W}/\beta$, and the unit
normalisation $\kappa_g = 1$ in natural rapidity units, all of the
following hold at every tested $(\nmax, N_t) \in \{1, 2, 3\} \times \{1, 11, 21\}$:
\begin{enumerate}\setlength{\itemsep}{2pt}
\item[(a)] \emph{BW-$\alpha$ period closure (geometric).}\ \ %
For every $O \in B(\mathcal{K}_{W_L})$,
$\sigma_{2\pi}^{L, \alpha}(O) = O$ bit-exactly, with periodicity residual
$\le 4 \times 10^{-16}$ (machine precision).
\item[(b)] \emph{BW-$\gamma$ Tomita period closure.}\ \ %
On the Krein-GNS Hilbert--Schmidt space
$\mathcal{K}_{\mathrm{GNS}} = M_{\dim W_L}(\mathbb{C})$ with modular
operator $\Delta_L$ and modular Hamiltonian $K_L^{\mathrm{TT}} = -\log \Delta_L$,
$\sigma_{2\pi}^{L, \mathrm{TT}}(a) = a$ bit-exactly at the same
$\le 4 \times 10^{-16}$ residual.
\item[(c)] \emph{Flow conjugacy at general $t$.}\ \ %
$\sigma_t^{L, \mathrm{TT}}(a) = \sigma_{-t}^{L, \alpha}(a)$ bit-exactly
for $t \in \{1, \pi, 2\pi\}$ on five sampled multipliers, residual
$\le 4 \times 10^{-16}$.
\item[(d)] \emph{Six-witness collapse.}\ \ %
For all six instantiations \{\textsc{bw}, $\textsc{hh}_{M=1}$,
$\textsc{hh}_{M=2}$, $\textsc{sew}_{M=1}$, $\textsc{unruh}_{a=1}$,
$\textsc{unruh}_{a=2}\}$, the cross-witness consistency residual is
exactly $0.0$.
\end{enumerate}
\end{theorem}

\begin{proof}[Proof sketch]
$K_L^{\alpha, W}$ is diagonal with integer eigenvalues
$\mathrm{two}\_m_j$ (inherited from Definition~\ref{def:K_alpha_W} on
the spatial side, with the temporal slot acting as $I_{N_t}$),
so $e^{i \cdot 2\pi \cdot K_L^{\alpha, W}} = I$ identically — proving
(a).  Statement (b) follows from
$\sigma_t^{L, \mathrm{TT}}(a) = e^{-i t K_L^{\alpha, W}} \, a \, e^{+i t K_L^{\alpha, W}}$
(polar decomposition + $\rho_W^L = e^{-K_L^{\alpha, W}}/Z$ identity)
combined with the integer-spectrum property of $K_L^{\alpha, W}$.
Statement (c) is the same operator identity at general $t$.  Statement
(d) follows from the $\beta$-independence of $\rho_W^L$ under the BW
choice $H_{\mathrm{local}} = K_L^{\alpha, W}/\beta$:\ the witnesses
parameterise different physical $\kappa_g$, but the modular
construction does not depend on $\kappa_g$ at the algebra-action
level.  Full proofs mirror Theorems~\ref{thm:bw_alpha},
\ref{thm:bw_gamma}, \ref{thm:unified}, and
Corollary~\ref{cor:single_construction} of the Riemannian closure;
the numerical verification at every panel cell is documented in
\texttt{debug/data/l2\_e\_modular\_hamiltonian\_lorentzian.json} via
\texttt{geovac/modular\_hamiltonian\_lorentzian.py}.
\end{proof}

The headline reading:\ the four-witness Wick-rotation theorem
(Hartle--Hawking + Sewell + Bisognano--Wichmann + Unruh) is lifted
from \emph{structural correspondence at the metric-functional level}
(Sprint TD~Track~4 + Sprint Unruh-pendant) to \emph{literal
identification at the operator-system level (Lorentzian, finite
cutoff)} at every tested Krein cutoff.  This is the framework's
first operator-system-level literal identification of a Lorentzian
QFT theorem at finite cutoff.

\subsection{Riemannian-limit identification (bit-identical at $N_t = 1$)}
\label{sec:riemannian_limit_identification}

The Sprint L2 architecture is the genuine Lorentzian \emph{extension}
of the Riemannian closure of \S\S\ref{sec:setup}--\ref{sec:four_witness_collapse},
not a parallel-but-different construction:

\begin{proposition}[Riemannian-limit recovery at $N_t = 1$]
\label{prop:riemannian_limit_l2}
At $N_t = 1$, every Sprint L2 construction reduces bit-identically
(Frobenius residual exactly $0.0$, not merely machine-precision)
to the corresponding Riemannian-side object of this paper:
\begin{itemize}\setlength{\itemsep}{2pt}
\item $P_{W_L}|_{N_t=1} = P_W^{\mathrm{spatial}}$ (Definition~\ref{def:wedge});
\item $K_L^\alpha|_{N_t=1} = K_\alpha$ (Definition~\ref{def:K_alpha_W});
\item $K_L^{\alpha, W}|_{N_t=1} = K_\alpha^W$;
\item $\rho_W^L|_{N_t=1} = \rho_W$ (\S\ref{sec:kms_state});
\item $\Delta_L|_{N_t=1} = \Delta$ (\S\ref{sec:tt_construction});
\item $K_L^{\mathrm{TT}}|_{N_t=1} = K_{\mathrm{TT}}$ (\S\ref{sec:tt_construction}).
\end{itemize}
The bit-identical reduction is structural:\ at $N_t = 1$ the
temporal-derivative piece $\partial_t$ vanishes on the singleton
grid, and the Kronecker product with $I_1 = 1$ is the identity
operation.  Every Riemannian-side bit-exact property (Theorems~\ref{thm:bw_alpha},
\ref{thm:bw_gamma}, \ref{thm:unified}) is preserved verbatim at
$N_t = 1$.
\end{proposition}

\emph{Verification.}\ \ Bit-identical at $n_{\max} \in \{1, 2, 3\}$,
documented in \texttt{tests/test\_modular\_hamiltonian\_lorentzian.py}.

This Riemannian-limit identification is the load-bearing falsifier
that ensures Sprint L2-E is the Lorentzian extension of the present
paper, not a parallel construction.  The temporal slot adds genuine
new content at $N_t > 1$ (the temporal-derivative refinement of
\S\ref{sec:h_local_signature_independent} below) that is invisible at
$N_t = 1$.

\subsection{The $H_{\mathrm{local}}$ verdict at $(3, 1)$ —
signature-independent at Riemannian limit}
\label{sec:h_local_signature_independent}

This subsection states the substantive structural finding of
Sprint L2-E.  Open question O3 of \S\ref{sec:conclusion} (the
spectral-action vs.\ modular-Hamiltonian generator distinction
flagged in \S\ref{sec:derived_hamiltonian} as the
``derived-Hamiltonian finding'') is refined at finite cutoff:

\begin{theorem}[$H_{\mathrm{local}}$ signature-independence at Riemannian limit]
\label{thm:h_local_sig_independent}
At $N_t = 1$ (Riemannian limit on the Krein space at $(3, 1)$),
\[
\bigl\| H_{\mathrm{local}}|_{(3, 1), N_t = 1} - D_L^W|_{N_t = 1} \bigr\|_F
\;=\;
\bigl\| H_{\mathrm{local}}|_{(3, 0)} - D_{\mathrm{GV}}^W \bigr\|_F
\]
bit-exact at every tested $n_{\max} \in \{1, 2, 3\}$:\ values
$2.1332$, $6.5275$, $13.854$ respectively, matching the Riemannian
residual of the present paper at the same $n_{\max}$.  At $N_t > 1$,
the Lorentzian residual is refined upward by the temporal-derivative
contribution to $D_L^W$, growing as $O(\sqrt{N_t})$ but with the
Riemannian baseline preserved.
\end{theorem}

\emph{Mechanism.}\ \ At $N_t = 1$, $D_L = i \cdot D_{\mathrm{GV}}$
exactly (the temporal derivative $\partial_t$ vanishes on the
singleton grid).  Wedge restriction gives $D_L^W = i \cdot D_{\mathrm{GV}}^W$.
The norm $\|H_{\mathrm{local}} - i \cdot D_{\mathrm{GV}}^W\|_F$ equals
$\|H_{\mathrm{local}} - D_{\mathrm{GV}}^W\|_F$ when $H_{\mathrm{local}}$
is real (which it is — $K_L^{\alpha, W}/\beta$ has real integer-rational
diagonal entries).  So the residuals are bit-identical.

\emph{Structural reading.}\ \ The \S\ref{sec:derived_hamiltonian} /
O3 distinction between the spectral-action Dirac $D_W$ and the
modular-Hamiltonian generator $K_\alpha^W/\beta$ on the BW vacuum is
\textbf{not a peculiarity of the round-$\sthree$ Riemannian
truncation}.  It persists bit-exactly at signature $(3, 1)$ at the
Riemannian limit, and is refined upward at $N_t > 1$ by the
temporal-derivative content of $D_L$.  The distinction is therefore
a \textbf{deeper structural feature} of the framework's modular
content, independent of signature at the Riemannian-limit baseline.

The temporal-derivative refinement at $N_t > 1$ is a Lorentzian-side
enrichment of the Riemannian finding:\ at $(3, 1)$, $D_L$ has more
content than $D_{\mathrm{GV}}$ (it includes the temporal derivative),
so the gap between the spectral-action object and the
modular-Hamiltonian generator widens with temporal cutoff.  This is
not a sign that \S\ref{sec:derived_hamiltonian} was wrong;\ it is a
sign that the Lorentzian extension reveals \emph{additional}
structure in the gap.

This finding refines Paper~42 O3 from ``open question about
signature-dependence'' to ``signature-INDEPENDENT structural
distinction at the Riemannian limit, refined by temporal-derivative
content at $N_t > 1$.''  The \emph{deeper} question (why is the
spectral-action Dirac not the modular generator on the BW vacuum,
structurally?) remains the residual O3 open question.

\subsection{Honest scope at finite cutoff}
\label{sec:lorentzian_closure_scope}

The closure of Theorem~\ref{thm:krein_four_witness} is at finite
$n_{\max}$ and finite $N_t$ — \textbf{not} in the continuum
(GH-limit, Lorentzian propinquity) limit.  Explicitly:

\begin{itemize}\setlength{\itemsep}{2pt}
\item The Latr\'emoli\`ere
propinquity~\cite{latremoliere2018, latremoliere_metric_st_2017} used
in Papers~38--40 for the Riemannian-side GH-convergence theorem is
purely Riemannian as of May~2026.  Subsequent
literature (Toyota~\cite{toyota2023}, Farsi--Latr\'emoli\`ere
2024/2025~\cite{farsi_latremoliere2024, farsi_latremoliere2025},
Hekkelman--McDonald 2024 a/b~\cite{hekkelman_mcdonald2024,
hekkelman_mcdonald2024b}) covers only Riemannian / Hilbert-space
structures.  No published Lorentzian propinquity construction exists.

\item Constructing a Lorentzian propinquity is the named Sprint L3
target (6--12 months of original NCG mathematics).  The
Nieuviarts~2025~\cite{nieuviarts2025a, nieuviarts2025b_proceedings}
twist-morphism may shorten this budget if the construction applies
to GeoVac $\sthree = \mathrm{SU}(2)$ (odd-dimension caveat that
Sprint L0 audit flagged for a separate L2-Nieuviarts-scoping pass).

\item The Sprint L2-D structural finding (BBB universal axiom
$\chi D = -D \chi$ fails on truthful Camporesi--Higuchi
$D_{\mathrm{GV}}$) is the load-bearing scope finding:\ the Sprint L2-E
construction uses truthful $D_{\mathrm{GV}}$ (which preserves the
Riemannian-limit recovery of Proposition~\ref{prop:riemannian_limit_l2}),
not a full BBB indefinite spectral triple in the strict Sec~5(v)
sense.  The wedge-modular construction is $K_L^{\alpha, W}$-driven,
not $D_L$-driven, so the period closure is $D_L$-independent at the
operator-action level.  Paper~32~\S~VIII.D~\cite{paper32} documents
the three resolutions (R1 accept, R2 use offdiag CH, R3 redefine
$\gamma^5$);\ Theorem~\ref{thm:krein_four_witness} uses R1 (truthful
$D_{\mathrm{GV}}$).
\end{itemize}

The verdict \textbf{closed-at-finite-cutoff} is therefore an honest
finite-cutoff statement.  Theorem~\ref{thm:krein_four_witness} is
the framework's first operator-system literal identification of a
Lorentzian QFT theorem;\ the Lorentzian-propinquity rate-level
extension is the named open Sprint L3 target.

\subsection{References and pointer to standalone writeup}
\label{sec:lorentzian_closure_references}

The construction recipe is van~den~Dungen 2016 Proposition~4.1~%
(arXiv:1505.01939):\ Wick-rotate $g \to g_r$ (positive-definite),
construct the standard Riemannian Dirac $\not{\!\!D}_{g_r}$, and set
$D_L = i^t \cdot \not{\!\!D}_{g_r}$ as the Krein-self-adjoint
Lorentzian Dirac with $\mathcal{J} = \gamma^0$.  At $(s, t) = (3, 1)$
this gives $i^t = i$ and the construction of
\S\ref{sec:lorentzian_closure_architecture}.

The Bizi--Brouder--Besnard 2018 $(m, n) \in (\mathbb{Z}/8)^2$
classification~\cite{bizi_brouder_besnard2018} provides the
signature-dependent Connes axiom signs:\ the $(s, t) = (3, 1)$
West-coast entry sits at $(m, n) = (4, 6)$ via BBB Table~3,
verified directly from the BBB PDF in Sprint L2-D (Paper~32~\S~VIII.D).

The full math.OA-style writeup of the Lorentzian extension is
outlined in
\texttt{papers/standalone/paper\_43\_lorentzian\_extension\_outline.md}
(companion document).  When drafted as a Zenodo deposit, Paper~43
should \textbf{not supersede} the present paper:\ Paper~42 is the
Riemannian closure of the four-witness theorem, Paper~43 is the
Lorentzian extension, and together they constitute the
four-witness Wick-rotation arc at the operator-system level.

The sprint memos in \texttt{debug/l2\_\{a,b,c,d,e,paper\_edits\_round1\}\_*.md}
and the synthesis memo \texttt{debug/sprint\_l2\_synthesis\_memo.md}
document the full Sprint L2 architecture and the substantive
structural findings; the computational drivers and verification data
sit in \texttt{geovac/\{krein\_space\_construction.py,
lorentzian\_dirac.py, connes\_axiom\_audit\_31.py,
modular\_hamiltonian\_lorentzian.py\}} and the corresponding
\texttt{tests/test\_*.py} files.

% =====================================================================
\appendix
\section{Computational verification protocol}
\label{app:computational_verification}
% =====================================================================

The numerical verification of the period closures of
Theorems~\ref{thm:bw_alpha} and~\ref{thm:bw_gamma} and of the
cross-witness collapse of Table~\ref{tab:cross_witness} is
implemented in the supporting Python module
\verb|geovac/modular_hamiltonian.py| of the present author's framework
codebase ($\sim$1550 lines), with test coverage in
\verb|tests/test_modular_hamiltonian.py| ($\sim$960 lines, 67 tests:\
63 fast plus 4 slow). The structure is briefly described here.

\subsection{Module API}\label{sec:module_api}

The principal classes are:
\begin{itemize}\setlength{\itemsep}{2pt}
\item \texttt{HemisphericWedge(axis="hopf")}:\ implements the projection
$P_W$ of Definition~\ref{def:wedge} via the equatorial-reflection
involution $R_{\mathrm{polar}}$ on the full-Dirac spinor basis.
\item \texttt{TomitaConjugation(U)}:\ implements the antilinear
$J_{\TT}$ via a stored unitary $U$ with action
$J_{\TT}(\psi) = U \cdot \overline{\psi}$;\ exposes
\texttt{verify\_J\_squared\_positive\_identity} and
\texttt{verify\_J\_squared\_negative\_identity\_DOES\_NOT\_HOLD} as
explicit sanity checks distinguishing $J_{\TT}$ from
$J_{\mathrm{GV}}$.
\item \texttt{KMSState(beta, H)}:\ implements the wedge-restricted
Gibbs state $\omega_{\beta, H}$ of~\eqref{eq:gibbs_state_def},
including the canonical analytic-continuation verification of the
KMS condition at the expectation level.
\item \texttt{ModularHamiltonian(wedge, basis, D, kappa\_g)}:\
implements the BW-$\alpha$ construction with $K = K_\alpha^W$, the
real-time flow $\sigma_t^{\alpha}$, the period-closure verifier
\texttt{verify\_modular\_periodicity}, and accessors to the
Tomita--Takesaki structure (\texttt{tomita\_structure},
\texttt{k\_alpha\_geometric}, \texttt{k\_tomita},
\texttt{compare\_alpha\_vs\_tomita}).
\item \texttt{TomitaModularStructure(rho)}:\ implements the
BW-$\gamma$ Tomita construction on the GNS Hilbert--Schmidt space,
with $\Delta = (\rho^{-1})^T \otimes \rho$, $K_{\TT} = -\log \Delta$,
and the modular flow $\sigma_t^{\TT}(a) = \rho^{it} a \rho^{-it}$ on
the algebra.
\end{itemize}
Witness factories \texttt{for\_bisognano\_wichmann},
\texttt{for\_hartle\_hawking}, \texttt{for\_sewell},
\texttt{for\_unruh} instantiate
\texttt{ModularHamiltonian} with witness-specific $\kappa_g$ values:\
BW canonical ($\kappa_g = 1$), HH ($\kappa_g = 1/(4M)$),
Unruh ($\kappa_g = a$). Each factory returns a
\texttt{ModularHamiltonian} object whose period-closure verifier
returns \texttt{STRONG\_IDENTIFICATION} at every tested $\nmax$.

\subsection{Test coverage}\label{sec:test_coverage}

The test suite \verb|tests/test_modular_hamiltonian.py| verifies the
following:\ wedge projector idempotency and Hermiticity ($P_W^2 =
P_W$, $P_W^* = P_W$); Tomita conjugation properties ($J_{\TT}^2 = +I$,
$J_{\TT}^2 \ne -I$, antilinearity, operator-action consistency); KMS
state normalisation, partition-function positivity, real expectation
on Hermitian operators; modular periodicity $\sigma_{2\pi}(O) = O$ at
$\nmax \in \{2, 3, 4\}$ for the first five non-identity multipliers in
\verb|FullDiracTruncatedOperatorSystem|; cross-witness collapse at
the operator-action level; KMS expectation-level condition for sample
operator pairs $(A, B)$; propinquity rate cross-check
(residual $\ll \gamma_{\nmax}$); operator-system non-leakage
($\sigma_{2\pi}(M)$ stays in $\Op_{\nmax}$); modular flow group law
$\sigma_t \cdot \sigma_s = \sigma_{t+s}$ and identity at $t = 0$. The
slow tests verify the closure at $\nmax = 5$ ($\sim$80 seconds wall
time per slow test).

\subsection{Wall time and reproducibility}\label{sec:wall_time}

Per-$\nmax$ wall time on a standard laptop (Intel i7, Python 3.11,
NumPy 1.26):\ $\sim 0.4$ s at $\nmax = 2$, $\sim 3.5$ s at $\nmax = 3$,
$\sim 45$ s at $\nmax = 4$, $\sim 260$ s at $\nmax = 5$ for the
BW-$\alpha$ construction;\ $\sim 0.2$ s, $\sim 3.7$ s, $\sim 45$ s,
$\sim 545$ s for the BW-$\gamma$ construction (the $\nmax = 5$ cost
dominated by the $4900 \times 4900$ GNS-Hilbert--Schmidt
diagonalisation). Beyond $\nmax = 5$ the GNS construction becomes
expensive; the structural reading of the closure is unchanged at
higher $\nmax$ since the integer-spectrum mechanism holds at any
cutoff. Computational drivers and structured output JSON are
available in the present author's framework documentation under
\verb|debug/l1_modular_hamiltonian_compute.py| and
\verb|debug/l1_tighten_tomita_compute.py| respectively.

% =====================================================================
\section{Cross-references to the GeoVac framework}
\label{app:framework_crossref}
% =====================================================================

The present paper is part of an integrated NCG-research programme
developed by the present author. We briefly note the principal
cross-references for the reader's orientation.

\paragraph{Paper~38 (\cite{paper38}, propinquity convergence on
$\sthree$).}
The Latr\'emoli\`ere quantum Gromov--Hausdorff propinquity convergence
of the truncated Camporesi--Higuchi spectral triples $\Tcal_{\nmax}$
to the continuum $\Tcal_{\sthree}$ is established in Paper~38 with
explicit rate~\eqref{eq:paper38_main}. The L1' / L2 / L3 / L4 / L5
lemmas of Paper~38 supply the metric-side analytic infrastructure
that the present paper builds on. The propinquity rate constant
$4/\pi = \Vol(S^2)/\pi^2$ identifies as the same Hopf-base measure
factor that drives the modular period closure of the present
Theorems~\ref{thm:bw_alpha}--\ref{thm:bw_gamma}.

\paragraph{Paper~40 (\cite{paper40_unified}, universal $4/\pi$ on
compact Lie groups).}
The Paper~38 propinquity convergence extends to all compact connected
Lie groups with bi-invariant metric, with the rate constant $4/\pi$
universal across the class via the Plancherel-weight $\times$
Vandermonde-Jacobian cancellation theorem (Paper~40~\S~3.2). The
modular structure of the present paper plausibly extends to the
higher-rank cases (Section~\ref{sec:conclusion} O4); the integer
spectrum of the boost-class generator $K_\alpha^W$ on the spinor
harmonic basis is a property of $\SU(2)$ (where the half-integer
weights double to integers), and the higher-rank analog requires
verifying that the relevant Cartan generator along the wedge axis
still has the right spectral property.

\paragraph{Paper~32 \S VIII (master Mellin engine case-exhaustion
theorem).}
The case-exhaustion theorem~\cite[Thm.~thm:pi\_source\_case\_exhaustion]{paper32}
of Paper~32 \S VIII (in the present author's framework documentation)
states that every $\pi$ in any finite chain of Paper~34 projections
engages one of three master Mellin engine sub-mechanisms:\ M1 (Hopf-base
measure $\Vol(S^2)/\pi^2$ at operator order $k = 0$), M2 (Seeley--DeWitt
$\sqrt{\pi}\, \Q$ at $k = 2$), or M3 (vertex-parity Hurwitz / Catalan
$G$ / Dirichlet $L$ at $k = 1$). The $2\pi$ in the modular period of
the present paper is the M1 signature at the operator-system level.
The Bisognano--Wichmann reading of M1 on temporal compactifications
(\cite[Rem.~rem:bisognano\_wichmann\_reading]{paper32}) is exactly the
structural correspondence that the present paper lifts to literal
identification.

\paragraph{Paper~34 \S III.27 (Wick rotation projection).}
The Wick rotation projection in the present author's projection
taxonomy~\cite{paper34} is the structural mechanism through which the
framework's Euclidean spectral content is read as a Lorentzian
physical observable. The Sprint L0 audit
(\texttt{debug/lorentzian\_l0\_audit\_memo.md}) classifies this
projection in the WICK\_MAP\_FREE bucket at the metric-functional
level. The present paper extends the projection's status to literal
identification at the operator-system level (Riemannian), reducing
the framework's reliance on the published Wick-rotation chain for the
period-closure identity itself.

\paragraph{Paper~24 (\cite{paper24}, Coulomb / HO asymmetry).}
The cross-manifold blocker named in Paper~24 \S V (Riemannian $\sthree$
$\otimes$ holomorphic-sector $S^5$ Bargmann--Segal triple) carries
through to the modular construction:\ the present paper does not
address the cross-manifold modular structure, and the natural
extension would require a mixed Riemannian / Hardy-sector
tensor-product propinquity that does not yet exist in the literature.

\paragraph{Paper~2 (\cite{paper2}, fine-structure-constant
observation).}
The present paper does not impinge on Paper~2's status as a
conjectural observation. The $\nmax = 3$ cutoff of
Table~\ref{tab:bw_alpha_residuals} happens to coincide with Paper~2's
selection-principle truncation ($\dim \Hilb_3 = 40 = g_3^{\mathrm{Dirac}}
= \Delta^{-1}$); this is a structural coincidence at the Hilbert-space
dimension level, not a derivation of Paper~2's combination rule
$K = \pi(B + F - \Delta)$. The combination rule remains a numerical
observation without first-principles derivation.

% =====================================================================
\section*{Acknowledgments}
% =====================================================================

The author thanks Alain Connes, Carlo Rovelli, Walter van Suijlekom,
Fran\c{c}ois Latr\'emoli\`ere, Edward Hekkelman, Edward McDonald,
Lukas Leimbach, Matilde Marcolli, Horacio Casini, Marina Huerta,
Robert Myers, and Rainer Verch for the foundational frameworks
on which this paper builds; remaining errors are the author's own.
This work was developed in an AI-augmented agentic workflow with
Anthropic's Claude:\ implementation of the supporting computational
modules, drafting of the internal proof memos used as supporting
documentation, and bibliographic collation were performed
collaboratively. All theorem statements, design choices, and
editorial decisions are the author's responsibility. The author
thanks the open-source NCG community (notably the SageMath, sympy,
mpmath, and scipy teams) for the computational infrastructure used in
the supporting numerical verifications.

% =====================================================================
\bibliographystyle{plainnat}
\begin{thebibliography}{99}

\bibitem[Bisognano \& Wichmann(1975)]{bisognano_wichmann1975}
J.~J.~Bisognano and E.~H.~Wichmann,
\newblock \emph{On the duality condition for a Hermitian scalar field},
\newblock J.~Math.~Phys. \textbf{16} (1975), 985--1007.

\bibitem[Bisognano \& Wichmann(1976)]{bisognano_wichmann1976}
J.~J.~Bisognano and E.~H.~Wichmann,
\newblock \emph{On the duality condition for quantum fields},
\newblock J.~Math.~Phys. \textbf{17} (1976), 303--321.

\bibitem[Bizi, Brouder \& Besnard(2018)]{bizi_brouder_besnard2018}
N.~Bizi, C.~Brouder, and F.~Besnard,
\newblock \emph{Space and time dimensions of algebras with
application to Lorentzian noncommutative geometry and quantum
electrodynamics},
\newblock J.~Math.~Phys. \textbf{59} (2018), 062303;
\newblock arXiv:1611.07062.

\bibitem[Camporesi \& Higuchi(1996)]{camporesi_higuchi1996}
R.~Camporesi and A.~Higuchi,
\newblock \emph{On the eigenfunctions of the Dirac operator on
spheres and real hyperbolic spaces},
\newblock J.~Geom.~Phys. \textbf{20} (1996), 1--18.

\bibitem[Casini \& Huerta(2009)]{casini_huerta2009}
H.~Casini and M.~Huerta,
\newblock \emph{Entanglement entropy in free quantum field theory},
\newblock J.~Phys. A: Math.~Theor. \textbf{42} (2009), 504007;
\newblock arXiv:0905.2562.

\bibitem[Casini, Huerta \& Myers(2011)]{casini_huerta_myers2011}
H.~Casini, M.~Huerta, and R.~C.~Myers,
\newblock \emph{Towards a derivation of holographic entanglement
entropy},
\newblock JHEP \textbf{05} (2011), 036;
\newblock arXiv:1102.0440.

\bibitem[Chamseddine \& Connes(2010)]{chamseddine_connes2010}
A.~H.~Chamseddine and A.~Connes,
\newblock \emph{Noncommutative geometry as a framework for
unification of all fundamental interactions including gravity. I},
\newblock Fortsch.~Phys. \textbf{58} (2010), 553--600.

\bibitem[Connes(1995)]{connes1995}
A.~Connes,
\newblock \emph{Noncommutative Geometry},
\newblock Academic Press, 1995.

\bibitem[Connes \& Rovelli(1994)]{connes_rovelli1994}
A.~Connes and C.~Rovelli,
\newblock \emph{Von Neumann algebra automorphisms and
time-thermodynamics relation in general covariant quantum theories},
\newblock Class.~Quantum Grav. \textbf{11} (1994), 2899--2918;
\newblock arXiv:gr-qc/9406019.

\bibitem[Connes \& van~Suijlekom(2021)]{connes_vs2021}
A.~Connes and W.~D.~van~Suijlekom,
\newblock \emph{Spectral truncations in noncommutative geometry and
operator systems},
\newblock Comm.~Math.~Phys. \textbf{383} (2021), 2021--2067;
\newblock arXiv:2004.14115.

\bibitem[Farsi \& Latr\'emoli\`ere(2024)]{farsi_latremoliere2024}
C.~Farsi and F.~Latr\'emoli\`ere,
\newblock \emph{Collapse in noncommutative geometry and spectral
continuity},
\newblock arXiv:2404.00240, 2024.

\bibitem[Farsi \& Latr\'emoli\`ere(2025)]{farsi_latremoliere2025}
C.~Farsi and F.~Latr\'emoli\`ere,
\newblock \emph{Continuity for the spectral propinquity of the
Dirac operators associated with an analytic path of Riemannian
metrics},
\newblock arXiv:2504.11715, 2025.

\bibitem[Hartle \& Hawking(1976)]{hartle_hawking1976}
J.~B.~Hartle and S.~W.~Hawking,
\newblock \emph{Path-integral derivation of black-hole radiance},
\newblock Phys.~Rev. D \textbf{13} (1976), 2188--2203.

\bibitem[Hekkelman(2022)]{hekkelman2022}
E.~Hekkelman,
\newblock \emph{Truncations of the circle and Connes' geometric
formula},
\newblock M.Sc.\ thesis, Radboud University, 2022;
\newblock arXiv:2206.13744.

\bibitem[Hekkelman \& McDonald(2024)]{hekkelman_mcdonald2024}
E.~Hekkelman and E.~McDonald,
\newblock \emph{Spectral truncations of $T^d$ and quantum metric
geometry},
\newblock J.~Noncommut.~Geom., to appear (2024);
\newblock arXiv:2403.18619.

\bibitem[Hekkelman \& McDonald(2024b)]{hekkelman_mcdonald2024b}
E.~Hekkelman and E.~McDonald,
\newblock \emph{A noncommutative integral on spectrally truncated
spectral triples, and a link with quantum ergodicity},
\newblock J.~Funct.~Anal., to appear (2025);
\newblock arXiv:2412.00628.

\bibitem[Hekkelman, McDonald \& van~Suijlekom(2024)]{ucp_maps_2024}
E.~Hekkelman, E.~McDonald, and W.~D.~van~Suijlekom,
\newblock \emph{UCP maps in spectral truncations of compact metric
spaces},
\newblock arXiv:2410.15454, 2024.

\bibitem[Latr\'emoli\`ere(2018)]{latremoliere2018}
F.~Latr\'emoli\`ere,
\newblock \emph{The Gromov--Hausdorff propinquity},
\newblock Trans.~Amer.~Math.~Soc. \textbf{370} (2018), 365--411.

\bibitem[Latr\'emoli\`ere(2017)]{latremoliere_metric_st_2017}
F.~Latr\'emoli\`ere,
\newblock \emph{The Gromov--Hausdorff propinquity for metric spectral
triples},
\newblock Adv.~Math. \textbf{415} (2023), Paper No.~108876, 88pp;
\newblock preprint arXiv:1811.10843, 2017.

\bibitem[Leimbach \& van~Suijlekom(2024)]{leimbach_vs2024}
L.~Leimbach and W.~D.~van~Suijlekom,
\newblock \emph{Gromov--Hausdorff convergence of spectral truncations
of the torus},
\newblock Adv.~Math. \textbf{439} (2024), Paper No.~109496.

\bibitem[Marcolli \& van~Suijlekom(2014)]{marcolli_vs2014}
M.~Marcolli and W.~D.~van~Suijlekom,
\newblock \emph{Gauge networks in noncommutative geometry},
\newblock J.~Geom.~Phys. \textbf{75} (2014), 71--91;
\newblock arXiv:1301.3480.

\bibitem[Nieuviarts(2024)]{nieuviarts2024}
A.~Nieuviarts,
\newblock \emph{From twisted spectral triples to pseudo-Riemannian
spectral triples in the almost commutative framework},
\newblock arXiv:2402.05839, 2024 (revised through v6, March 2026).

\bibitem[Nieuviarts(2025a)]{nieuviarts2025a}
A.~Nieuviarts,
\newblock \emph{Emergence of pseudo-Riemannian spectral triples
within the almost-commutative framework},
\newblock arXiv:2502.18105 (v3, May 2025).

\bibitem[Nieuviarts(2025b)]{nieuviarts2025b_proceedings}
A.~Nieuviarts,
\newblock \emph{Emergence of pseudo-Riemannian structures from
twisted spectral triples within the almost-commutative framework
(proceedings synthesis)},
\newblock arXiv:2512.15450 (v2, May 2026).

\bibitem[Sewell(1982)]{sewell1982}
G.~L.~Sewell,
\newblock \emph{Quantum fields on manifolds: PCT and gravitationally
induced thermal states},
\newblock Ann.~Phys. (NY) \textbf{141} (1982), 201--224.

\bibitem[Strohmaier(2006)]{strohmaier2006}
A.~Strohmaier,
\newblock \emph{The Reeh--Schlieder property for quantum fields on
stationary spacetimes},
\newblock Comm.~Math.~Phys. \textbf{215} (2000), 105--118.

\bibitem[Takesaki(1970)]{takesaki1970}
M.~Takesaki,
\newblock \emph{Tomita's theory of modular Hilbert algebras and its
applications},
\newblock Lecture Notes in Mathematics 128, Springer, 1970.

\bibitem[Tomita(1967)]{tomita1967}
M.~Tomita,
\newblock \emph{Standard form of von Neumann algebras},
\newblock Lecture notes, RIMS Kokyuroku \textbf{14}, Kyoto University,
1967.

\bibitem[Toyota(2023)]{toyota2023}
M.~Toyota,
\newblock \emph{Quantum Gromov--Hausdorff convergence of spectral
truncations for groups with polynomial growth},
\newblock arXiv:2309.13469, 2023.

\bibitem[Unruh(1976)]{unruh1976}
W.~G.~Unruh,
\newblock \emph{Notes on black-hole evaporation},
\newblock Phys.~Rev. D \textbf{14} (1976), 870--892.

\bibitem[Verch(2001)]{verch2001}
R.~Verch,
\newblock \emph{A spin-statistics theorem for quantum fields on
curved spacetime manifolds in a generally covariant framework},
\newblock Comm.~Math.~Phys. \textbf{223} (2001), 261--288.

\bibitem[Zhu, Casini, \emph{et al.}(2020)]{zhu_casini2020}
J.~Zhu, H.~Casini, M.~Dalmonte, and P.~Hauke,
\newblock \emph{Lattice Bisognano--Wichmann modular Hamiltonian in
critical quantum spin chains},
\newblock SciPost Phys.~Core \textbf{2} (2020), 007;
\newblock arXiv:2003.00315.

\bibitem[Loutey, internal Paper~2]{paper2}
J.~Loutey,
\newblock \emph{Fine structure constant from the Hopf bundle on $S^3$:\
$K = \pi(B + F - \Delta)$ as a three-sector spectral coincidence},
\newblock GeoVac internal preprint, papers/observations (2026),
DOI via Zenodo.

\bibitem[Loutey, internal Paper~24]{paper24}
J.~Loutey,
\newblock \emph{Bargmann--Segal lattice for the 3D harmonic oscillator
and the Coulomb/HO structural asymmetry},
\newblock GeoVac internal preprint (2026), DOI via Zenodo.

\bibitem[Loutey, internal Paper~25]{paper25_hopf_gauge_structure}
J.~Loutey,
\newblock \emph{Hopf graph as a discrete lattice-gauge structure on
the Fock-projected $S^3$},
\newblock GeoVac internal preprint, papers/synthesis (2026),
DOI via Zenodo.

\bibitem[Loutey, internal Paper~32]{paper32}
J.~Loutey,
\newblock \emph{The GeoVac spectral triple, the master Mellin engine
case-exhaustion theorem, and the GH-convergence theorem on $\sthree$},
\newblock GeoVac internal preprint, papers/synthesis (2026),
DOI via Zenodo.

\bibitem[Loutey, internal Paper~34]{paper34}
J.~Loutey,
\newblock \emph{Projection taxonomy and empirical matches catalogue
for the GeoVac framework},
\newblock GeoVac internal preprint, papers/observations (2026),
DOI via Zenodo.

\bibitem[Loutey, internal Paper~38]{paper38}
J.~Loutey,
\newblock \emph{Latr\'emoli\`ere propinquity convergence of
truncated Camporesi--Higuchi spectral triples on $\sthree$},
\newblock GeoVac internal preprint, papers/standalone (2026),
DOI via Zenodo.

\bibitem[Loutey, internal Paper~39]{paper39}
J.~Loutey,
\newblock \emph{Tensor-product propinquity convergence for
distinct-focal-length truncated Camporesi--Higuchi spectral triples
on $\sthree$},
\newblock GeoVac internal preprint, papers/standalone (2026),
DOI via Zenodo.

\bibitem[Loutey, internal Paper~40]{paper40_unified}
J.~Loutey,
\newblock \emph{Latr\'emoli\`ere propinquity convergence of spectral
truncations of compact Lie groups with bi-invariant metric},
\newblock GeoVac internal preprint, papers/standalone (2026),
DOI via Zenodo.

\bibitem[Loutey, internal Paper~43]{paper43}
J.~Loutey,
\newblock \emph{Lorentzian extension of the four-witness
Wick-rotation theorem at finite cutoff},
\newblock GeoVac internal preprint, papers/standalone (2026),
DOI via Zenodo.

\end{thebibliography}

\end{document}
